\documentclass[11pt,oneside]{article}
% Standalone ProveIt research note.  The source intentionally has no external
% inputs so that the TeX/PDF pair is directly reproducible after extraction.
\usepackage[a4paper,margin=26mm,headheight=15pt]{geometry}
\usepackage[T1]{fontenc}
\usepackage[utf8]{inputenc}
\IfFileExists{libertinus.sty}{\usepackage{libertinus}}{\usepackage{lmodern}}
\usepackage{microtype}
\usepackage{amsmath,amssymb,amsthm,mathtools,mathrsfs,bm}
\usepackage{booktabs,longtable,tabularx,array,multirow}
\usepackage{graphicx}
\usepackage{xcolor}
\usepackage{enumitem}
\usepackage{fancyhdr}
\usepackage{xurl}
\usepackage{listings}
\usepackage[most]{tcolorbox}
\usepackage{hyperref}
\usepackage[nameinlink,noabbrev,capitalise]{cleveref}

\definecolor{linkblue}{RGB}{24,72,124}
\definecolor{deepgreen}{RGB}{23,102,69}
\definecolor{deepred}{RGB}{140,42,42}
\definecolor{deepgold}{RGB}{122,91,19}
\definecolor{shadegray}{RGB}{246,247,249}
\definecolor{rulegray}{RGB}{195,201,208}
\definecolor{palered}{RGB}{251,236,236}
\definecolor{palegold}{RGB}{251,245,230}
\definecolor{palegreen}{RGB}{233,245,243}
\definecolor{paleblue}{RGB}{236,243,251}

\hypersetup{
  hypertexnames=false,
  colorlinks=true,
  linkcolor=linkblue,
  citecolor=deepgreen,
  urlcolor=linkblue,
  pdftitle={Log-Periodic Transseries for the Inverse of a Real Fibonacci Continuation},
  pdfauthor={Vladimir Reshetnikov},
  pdfsubject={Asymptotic inversion of an exponential-oscillatory real continuation of the Fibonacci numbers},
  pdfkeywords={Fibonacci numbers, Binet formula, inverse function, transseries, log-periodic asymptotics, complex powers, Lagrange-Good inversion, oscillatory exponential asymptotics}
}

\pagestyle{fancy}
\fancyhf{}
\fancyhead[L]{\small Inverse Fibonacci transseries}
\fancyhead[R]{\small\nouppercase{\leftmark}}
\fancyfoot[C]{\thepage}
\renewcommand{\headrulewidth}{0.4pt}

\setcounter{tocdepth}{2}
\setcounter{secnumdepth}{3}
\setlist[itemize]{leftmargin=1.5em,itemsep=0.25ex,topsep=0.6ex}
\setlist[enumerate]{leftmargin=1.9em,itemsep=0.35ex,topsep=0.6ex}
\allowdisplaybreaks[2]
\setlength{\parindent}{1.2em}
\setlength{\parskip}{0.15em}

\theoremstyle{plain}
\newtheorem{theorem}{Theorem}[section]
\newtheorem{proposition}[theorem]{Proposition}
\newtheorem{lemma}[theorem]{Lemma}
\newtheorem{corollary}[theorem]{Corollary}
\theoremstyle{definition}
\newtheorem{definition}[theorem]{Definition}
\newtheorem{algorithm}[theorem]{Algorithm}
\newtheorem{example}[theorem]{Example}
\newtheorem{conjecture}[theorem]{Conjecture}
\theoremstyle{remark}
\newtheorem{remark}[theorem]{Remark}

\crefname{algorithm}{algorithm}{algorithms}
\Crefname{algorithm}{Algorithm}{Algorithms}
\crefname{conjecture}{conjecture}{conjectures}
\Crefname{conjecture}{Conjecture}{Conjectures}

\newtcolorbox{mainbox}[1][]{%
  enhanced,breakable,colback=paleblue,colframe=linkblue,
  boxrule=0.7pt,arc=1.2mm,left=2mm,right=2mm,top=1.4mm,bottom=1.4mm,
  title={Main result},fonttitle=\bfseries,#1}
\newtcolorbox{warningbox}[1][]{%
  enhanced,breakable,colback=palered,colframe=deepred,
  boxrule=0.6pt,arc=1.2mm,left=2mm,right=2mm,top=1.2mm,bottom=1.2mm,
  title={Scale warning},fonttitle=\bfseries,#1}
\newtcolorbox{methodbox}[1][]{%
  enhanced,breakable,colback=palegreen,colframe=deepgreen,
  boxrule=0.6pt,arc=1.2mm,left=2mm,right=2mm,top=1.2mm,bottom=1.2mm,
  title={Reusable method},fonttitle=\bfseries,#1}
\newtcolorbox{summarybox}[1][]{%
  enhanced,breakable,colback=palegold,colframe=deepgold,
  boxrule=0.6pt,arc=1.2mm,left=2mm,right=2mm,top=1.2mm,bottom=1.2mm,
  title={Summary},fonttitle=\bfseries,#1}

\lstdefinestyle{codestyle}{%
  basicstyle=\ttfamily\small,
  columns=fullflexible,
  frame=single,
  rulecolor=\color{rulegray},
  backgroundcolor=\color{shadegray},
  showstringspaces=false,
  keepspaces=true,
  breaklines=true,
  xleftmargin=0.4em,xrightmargin=0.4em,
  aboveskip=0.8em,belowskip=0.8em,
  keywordstyle=\color{linkblue}\bfseries,
  commentstyle=\color{deepgreen},
  stringstyle=\color{deepred}
}
\lstdefinestyle{wolframstyle}{style=codestyle,language=Mathematica}
\lstdefinestyle{pythonstyle}{style=codestyle,language=Python}

\newcommand{\e}{\mathrm{e}}
\newcommand{\ii}{\mathrm{i}}
\newcommand{\dd}{\mathop{}\!\mathrm{d}}
\newcommand{\R}{\mathbb{R}}
\newcommand{\C}{\mathbb{C}}
\newcommand{\N}{\mathbb{N}}
\newcommand{\Z}{\mathbb{Z}}
\newcommand{\Fcal}{\mathcal{F}}
\newcommand{\Hcal}{\mathcal{H}}
\newcommand{\Tcal}{\mathcal{T}}
\newcommand{\Mcal}{\mathcal{M}}
\newcommand{\Ocal}{\mathcal{O}}
\newcommand{\phig}{\varphi}
\newcommand{\coeff}[2]{\left[#1^{#2}\right]}
\newcommand{\mcoeff}[2]{\left[\bm{#1}^{\bm{#2}}\right]}
\newcommand{\abs}[1]{\left\lvert#1\right\rvert}
\newcommand{\norm}[1]{\left\lVert#1\right\rVert}
\newcommand{\avg}[1]{\left\langle#1\right\rangle}
\newcommand{\Real}{\operatorname{Re}}
\newcommand{\Imag}{\operatorname{Im}}
\newcommand{\arsinh}{\operatorname{arsinh}}
\newcommand{\arccsch}{\operatorname{arccsch}}
\newcommand{\sgn}{\operatorname{sgn}}
\newcommand{\Log}{\operatorname{Log}}
\DeclareMathOperator{\Res}{Res}
\DeclareMathOperator{\Trunc}{Trunc}

\begin{document}

\title{\textbf{Log-Periodic Transseries for the Inverse of a Real Fibonacci Continuation}\\[0.35em]
  \Large An Explicit Lagrange--Good Calculus for Exponential--Oscillatory Inversion}
\author{Vladimir Reshetnikov\\
  \small ProveIt research note}
\date{3 September 2026}
\maketitle

\begin{abstract}
Let
\[
 \Fcal(x)=\sqrt{\frac25}\,
 \sqrt{\cosh\!\bigl(2x\,\arccsch 2\bigr)-\cos(\pi x)},
 \qquad x\in\R.
\]
Because \(\arccsch 2=\log\phig\), where
\(\phig=(1+\sqrt5)/2\), this is the modulus of the real-parameter Binet
continuation and satisfies \(\Fcal(n)=\abs{F_n}\) at every integer.  The
function is even.  On \([2,\infty)\) it is strictly increasing from \(1\) to
\(\infty\), so it has an outer positive inverse \(\Psi(y)\); the corresponding
negative inverse is \(-\Psi(y)\).

This article obtains an exact all-orders expansion of these inverse branches as
\(y\to+\infty\).  Put
\[
 a=\log\phig,\qquad Y=\sqrt5\,y,\qquad
 t=Y^{-2},\qquad
 \theta=\frac{\pi}{a}\log Y,\qquad
 \rho=\frac{\pi}{2a}.
\]
Then, for all sufficiently large \(y\),
\[
 \Psi(y)=\frac{\log Y}{a}
 +\frac1{2a}\sum_{n\ge1}t^n Z_n(\theta),
\]
where \(Z_n\) is a real trigonometric polynomial with Fourier modes
\(-n,-n+2,\ldots,n\).  More strongly, with
\(\lambda=1-\ii\rho\), \(X=t^\lambda=t\e^{\ii\theta}\), and
\(A_{jk}=j\lambda+k\overline\lambda\), one has the convergent double
transseries
\[
 \Psi(y)=\frac{\log Y}{a}
 +\frac1{2a}\sum_{j+k\ge1}
 \frac{(-1)^{j+k+1}}{A_{jk}}
 \binom{A_{jk}}{j}\binom{A_{jk}}{k}
 X^j\overline X^{\,k}.
\]
The first six blocks are displayed explicitly, and triangular, fixed-point,
coefficient-extraction, and Newton algorithms are developed.  A residual
certificate converts every truncation into a rigorous numerical approximation.

The Fibonacci example is a specialization of a general inversion theorem.  If
\[
 Y=C\e^{\sigma x}
 H\!\left(c_1\e^{-\alpha_1x},\ldots,c_r\e^{-\alpha_rx}\right),
 \qquad H(0)=1,
\]
with \(H\) analytic and \(\Real\alpha_j>0\), then the inverse correction has
the explicit multivariate Lagrange--Good expansion
\[
 x(Y)=\frac1\sigma\log\frac YC
 +\frac1\sigma\sum_{\bm n\ne0}
 \left[-\frac1{A_{\bm n}}
   [\bm u^{\bm n}]H(\bm u)^{A_{\bm n}}\right]
 \prod_{j=1}^r\left(c_j(Y/C)^{-\alpha_j/\sigma}\right)^{n_j},
 \qquad
 A_{\bm n}=\sum_j n_j\frac{\alpha_j}{\sigma}.
\]
When a combined exponent \(A_{\bm n}\) vanishes, the displayed quotient has
a removable singularity: its regularized value is
\(-[\bm u^{\bm n}]\log H\).  Thus the common-factor inversion itself has no
genuine resonance pole.  The formula unifies real power corrections,
logarithmically periodic corrections from complex-conjugate exponents, and
quasiperiodic corrections from several incommensurable frequencies.  Product
and radical perturbations reduce to closed products of generalized binomial
coefficients.
\end{abstract}

\noindent\textbf{Keywords.}
Fibonacci numbers; Binet continuation; inverse function; transseries;
log-periodic asymptotics; complex powers; Lagrange--Good inversion;
multivariate Lagrange inversion; exponential--oscillatory asymptotics;
series reversion.

{\small\tableofcontents}
\clearpage

\section{The inverse problem and the answer in compact form}
\label{sec:intro}

The familiar Binet formula interpolates the Fibonacci sequence only after one
chooses what to do with the factor \((-1)^x\) away from the integers.  One
natural complex continuation traces a curve in the complex plane; taking its
modulus produces the real, nonnegative function considered here.  The
oscillation that remains after taking the modulus is exponentially small in
\(x\), but it is not negligible for a complete inverse expansion.  Upon
inversion, it becomes a hierarchy of powers of \(y^{-2}\) whose coefficients
are periodic functions of \(\log y\).

This is precisely the setting in which an ordinary Poincar\'e series with
constant coefficients is too coarse.  The correct scale is a
\emph{log-periodic transseries}: algebraic decay in the large variable,
multiplied by Fourier modes in its logarithm.  Equivalently, it is a series in
complex powers.  The two descriptions are identical because
\[
 y^{-\alpha-\ii\beta}=y^{-\alpha}\e^{-\ii\beta\log y}.
\]

The directory that motivated this note treats series and transseries as an
operational calculus of algebra, composition, and reversion
\cite{ProveItSeriesTransseries}.  The present example adds a different scale to
that program: a power scale with logarithmically oscillatory coefficient
blocks.

\begin{mainbox}[title={Outer inverse and closed all-orders expansion}]
Set
\begin{equation}
 a=\arccsch 2=\arsinh\frac12=\log\phig,
 \qquad \phig=\frac{1+\sqrt5}{2},
 \qquad \rho=\frac{\pi}{2a}.
 \label{eq:a-rho}
\end{equation}
The restriction \(\Fcal:[2,\infty)\to[1,\infty)\) is a bijection.  Denote its
inverse by \(\Psi\).  For \(y\to+\infty\), define
\begin{equation}
 Y=\sqrt5\,y,\qquad t=Y^{-2},\qquad
 \theta=\frac{\pi}{a}\log Y,
 \qquad \lambda=1-\ii\rho,
 \qquad X=t^\lambda=t\e^{\ii\theta}.
 \label{eq:scaled-variables}
\end{equation}
For \(j,k\ge0\), \(j+k\ge1\), put
\begin{equation}
 A_{jk}=j\lambda+k\overline\lambda
       =j+k-\ii\rho(j-k),
 \qquad
 d_{jk}=\frac{(-1)^{j+k+1}}{A_{jk}}
        \binom{A_{jk}}{j}\binom{A_{jk}}{k}.
 \label{eq:djk}
\end{equation}
Then
\begin{equation}
 \boxed{
 \Psi(y)=\frac{\log Y}{a}
 +\frac1{2a}\sum_{\substack{j,k\ge0\\j+k\ge1}}
 d_{jk}X^j\overline X^{\,k}.}
 \label{eq:main-double}
\end{equation}
The double series is absolutely convergent for all sufficiently large positive
\(y\).  Grouping by total degree gives
\begin{equation}
 \boxed{
 \Psi(y)=\frac{\log Y}{a}
 +\frac1{2a}\sum_{n\ge1}t^nZ_n(\theta),}
 \label{eq:main-block}
\end{equation}
where
\begin{equation}
 Z_n(\theta)=\sum_{j=0}^{n}
 \frac{(-1)^{n+1}}{n-\ii\rho(2j-n)}
 \binom{n-\ii\rho(2j-n)}{j}
 \binom{n-\ii\rho(2j-n)}{n-j}
 \e^{\ii(2j-n)\theta}.
 \label{eq:Zn-closed}
\end{equation}
The negative outer inverse is \(\Psi_-(y)=-\Psi(y)\).
\end{mainbox}

The first three nonzero blocks give the immediately usable expansion
\begin{align}
 \Psi(y)={}&\frac{\log(\sqrt5y)}a
 +\frac{\cos\theta}{5ay^2}
 -\frac{2+\cos(2\theta)+2\rho\sin(2\theta)}{50ay^4}
 \notag\\
 &+\frac1{250ay^6}\Bigl[
 (6-\rho^2)\cos\theta+5\rho\sin\theta
 +\Bigl(\frac23-3\rho^2\Bigr)\cos(3\theta)
 +3\rho\sin(3\theta)\Bigr]
 +O(y^{-8}),
 \label{eq:first-usable}
\end{align}
where \(\theta=(\pi/a)\log(\sqrt5y)\).  The oscillation has constant period
\(2a\) in \(\log Y\), or multiplicative period \(\e^{2a}=\phig^2\) in
\(Y\), at the level of a fixed Fourier mode.  Different blocks carry
harmonics of that base phase.

\begin{warningbox}[title={Why a constant-coefficient inverse series cannot be complete}]
The leading approximation \(\Psi(y)\sim a^{-1}\log(\sqrt5y)\) is correct, but
the first correction is
\[
 \frac{\cos\!\bigl((\pi/a)\log(\sqrt5y)\bigr)}{5ay^2}.
\]
Its coefficient never approaches a constant.  Consequently there is no
complete expansion of the form
\(a^{-1}\log(\sqrt5y)+\sum_{n\ge1}c_ny^{-2n}\) with constant \(c_n\).
One must retain either periodic coefficient functions or complex powers.
\end{warningbox}

The remainder of the article derives these formulas, proves convergence and
branch properties, gives several algorithms, and then develops the general
inversion theorem from which \eqref{eq:main-double} follows as a two-factor
special case.

\section{The real Binet--Fibonacci continuation}
\label{sec:function}

\subsection{Equivalent forms}

We begin by recording the algebraic identities that make the inversion
tractable.  Throughout,
\begin{equation}
 a=\arccsch 2=\arsinh\frac12=\log\phig>0.
 \label{eq:a-identities}
\end{equation}
The equality \(\arsinh(1/2)=\log\phig\) follows from
\[
 \arsinh\frac12=\log\left(\frac12+\sqrt{1+\frac14}\right)
 =\log\frac{1+\sqrt5}{2}.
\]
Squaring the defining formula gives
\begin{align}
 5\Fcal(x)^2
 &=2\cosh(2ax)-2\cos(\pi x) \notag\\
 &=\e^{2ax}+\e^{-2ax}-2\cos(\pi x) \label{eq:F-square-exp}\\
 &=\abs{\e^{ax}-\e^{-ax}\e^{\ii\pi x}}^2. \label{eq:F-modulus}
\end{align}
Thus one may equivalently write
\begin{equation}
 \Fcal(x)=\frac1{\sqrt5}
 \abs{\phig^x-\phig^{-x}\e^{\ii\pi x}}.
 \label{eq:F-complex-Binet}
\end{equation}
This is the modulus of the complex Binet continuation studied, in a broader
geometric setting, in \cite{OzvatanPashaev2017}.

\begin{proposition}[Interpolation, symmetry, and zeros]
\label{prop:basic-F}
The function \(\Fcal\) has the following properties.
\begin{enumerate}
\item \(\Fcal\) is continuous, nonnegative, and even on \(\R\).
\item \(\Fcal(x)=0\) if and only if \(x=0\).
\item For every integer \(n\),
\begin{equation}
 \Fcal(n)=\abs{F_n},
 \label{eq:integer-interpolation}
\end{equation}
where \(F_n\) is the usual two-sided Fibonacci sequence.  In particular,
\(\Fcal(n)=F_n\) for \(n\ge0\).
\item As \(x\to+\infty\),
\begin{equation}
 \Fcal(x)=\frac{\e^{ax}}{\sqrt5}
 \left(1+O(\e^{-2ax})\right).
 \label{eq:F-envelope}
\end{equation}
\end{enumerate}
\end{proposition}

\begin{proof}
Evenness and continuity are immediate from the original formula.  The modulus
representation \eqref{eq:F-modulus} can vanish only when
\(\e^{ax}=\e^{-ax}\e^{\ii\pi x}\).  Equality of moduli forces
\(\e^{ax}=\e^{-ax}\), hence \(x=0\); conversely \(\Fcal(0)=0\).

At an integer \(n\), \(\e^{\ii\pi n}=(-1)^n\), and Binet's formula gives
\[
 F_n=\frac{\phig^n-(-1)^n\phig^{-n}}{\sqrt5}.
\]
Equation \eqref{eq:F-complex-Binet} therefore yields
\(\Fcal(n)=\abs{F_n}\).  Finally, factor \(\e^{ax}\) from
\eqref{eq:F-complex-Binet}.  The remaining factor is the modulus of
\(1-\e^{-2ax+\ii\pi x}\), which equals \(1+O(\e^{-2ax})\).
\end{proof}

The modulus in \eqref{eq:F-complex-Binet} is essential.  Without it, the
complex continuation retains an exponentially decaying imaginary part.  With
it, the function is real but its square retains the oscillatory term
\(-2\cos(\pi x)\).  The inverse therefore remembers the phase \(\pi x\), even
though that phase appears only in an exponentially small relative correction
to the dominant envelope.

\subsection{The outer real inverse branches}

The function is not globally one-to-one on \([0,\infty)\); it has small
oscillations near the origin.  These oscillations are irrelevant for large
values, but the branch convention must be stated explicitly.

\begin{lemma}[A useful exact hyperbolic value]
\label{lem:sinh4a}
With \(a=\log\phig\),
\begin{equation}
 \sinh(2a)=\frac{\sqrt5}{2},\qquad
 \cosh(2a)=\frac32,
 \qquad
 \sinh(4a)=\frac{3\sqrt5}{2}.
 \label{eq:hyperbolic-values}
\end{equation}
Moreover,
\begin{equation}
 3\sqrt5\,a>\pi.
 \label{eq:key-inequality}
\end{equation}
\end{lemma}

\begin{proof}
Since \(\phig^2=(3+\sqrt5)/2\) and
\(\phig^{-2}=(3-\sqrt5)/2\), the first two identities follow by taking their
half-difference and half-sum.  The double-angle identity then gives
\(\sinh(4a)=2\sinh(2a)\cosh(2a)=3\sqrt5/2\).

For a completely elementary proof of the inequality, use
\((1+s^2)^{-1/2}\ge1-s^2/2\) for \(0\le s\le1/2\).  Hence
\[
 a=\int_0^{1/2}\frac{\dd s}{\sqrt{1+s^2}}
 \ge \frac12-\frac1{48}=\frac{23}{48}.
\]
Therefore \(3\sqrt5a\ge23\sqrt5/16\).  The classical bound
\(\pi<22/7\), together with
\(23\sqrt5/16>22/7\), proves \eqref{eq:key-inequality}.  The latter comparison
may be checked by squaring: \(5\cdot161^2>352^2\).
\end{proof}

\begin{theorem}[Monotonicity of the outer branch]
\label{thm:outer-monotone}
The restriction
\begin{equation}
 \Fcal:[2,\infty)\longrightarrow[1,\infty)
 \label{eq:outer-map}
\end{equation}
is a strictly increasing bijection.  Consequently there is a unique outer
positive inverse
\begin{equation}
 \Psi:[1,\infty)\longrightarrow[2,\infty),
 \qquad \Fcal(\Psi(y))=y.
 \label{eq:Psi-definition}
\end{equation}
By evenness, \(\Psi_-(y)=-\Psi(y)\) is the outer negative inverse.
\end{theorem}

\begin{proof}
Let
\(G(x)=\cosh(2ax)-\cos(\pi x)\).  Since \(\Fcal\) is a positive constant
times \(\sqrt{G}\), its derivative has the sign of
\[
 G'(x)=2a\sinh(2ax)+\pi\sin(\pi x).
\]
For \(x\ge2\), \Cref{lem:sinh4a} gives
\[
 2a\sinh(2ax)\ge2a\sinh(4a)=3\sqrt5a>\pi,
\]
whereas \(\pi\sin(\pi x)\ge-\pi\).  Thus \(G'(x)>0\) and hence
\(\Fcal'(x)>0\) on \([2,\infty)\).  Also \(\Fcal(2)=F_2=1\), while
\(\Fcal(x)\to\infty\) by \eqref{eq:F-envelope}.  The asserted bijection
follows.  Evenness supplies the negative branch.
\end{proof}

\begin{remark}[How many real inverse values?]
\label{rem:branch-count}
The compact interval \([-2,2]\) can contribute additional inverse values only
for bounded \(y\).  A simple rigorous bound is
\[
 \max_{\abs{x}\le2}\Fcal(x)
 \le \sqrt{\frac25\bigl(\cosh(4a)+1\bigr)}
 =\frac3{\sqrt5}.
\]
Hence, for \(y>3/\sqrt5\), the only real solutions of \(\Fcal(x)=y\) are
\(x=\pm\Psi(y)\).  Numerically, the actual final internal local maximum is much
smaller: it occurs near
\[
 x=1.1323790680,\qquad \Fcal(x)=1.0138598855,
\]
and is followed by a local minimum near
\[
 x=1.7192508860,\qquad \Fcal(x)=0.9111726845.
\]
These decimal values are descriptive, not needed in any proof below.
\end{remark}

\subsection{First-order inversion}

The envelope \eqref{eq:F-envelope} immediately suggests
\begin{equation}
 \Psi(y)\sim \frac1a\log(\sqrt5y).
 \label{eq:leading-inverse}
\end{equation}
The central problem is to calculate the complete correction.  The relative
perturbation in \(\Fcal(x)\) has size \(\e^{-2ax}\); after substituting the
leading inverse, this becomes \((\sqrt5y)^{-2}\).  The phase
\(\pi x\) becomes \((\pi/a)\log(\sqrt5y)\).  This dimensional argument already
predicts a series in \(y^{-2}\) with periodic coefficients in \(\log y\).  The
next section turns this prediction into an exact normalized equation.

\section{Exact normalization of the inverse equation}
\label{sec:normalization}

\subsection{The scalar implicit equation}

Let \(y\ge1\) and set \(x=\Psi(y)\).  Introduce
\begin{equation}
 Y=\sqrt5y,
 \qquad x_0=\frac{\log Y}{a},
 \qquad t=Y^{-2},
 \qquad \theta=\frac{\pi}{a}\log Y,
 \qquad \rho=\frac{\pi}{2a}.
 \label{eq:normalization-variables}
\end{equation}
Write the inverse in the anchored form
\begin{equation}
 x=x_0+\frac{z}{2a}.
 \label{eq:x-z-substitution}
\end{equation}
The unknown \(z=z(t,\theta)\) is small when \(t\) is small.

\begin{proposition}[Exact normalized equation]
\label{prop:normalized-equation}
The correction \(z\) satisfies
\begin{equation}
 \boxed{
 H(z;t,\theta):=
 \e^z+t^2\e^{-z}-2t\cos(\theta+\rho z)-1=0.}
 \label{eq:H-equation}
\end{equation}
Conversely, the solution of \eqref{eq:H-equation} that tends to zero with
\(t\) yields the outer inverse through \eqref{eq:x-z-substitution}.
\end{proposition}

\begin{proof}
The equation \(\Fcal(x)=y\), after squaring and multiplying by \(5\), is
\[
 Y^2=\e^{2ax}+\e^{-2ax}-2\cos(\pi x).
\]
Under \eqref{eq:x-z-substitution},
\[
 \e^{2ax}=Y^2\e^z,
 \qquad
 \e^{-2ax}=Y^{-2}\e^{-z},
 \qquad
 \pi x=\theta+\rho z.
\]
Division by \(Y^2\), with \(t=Y^{-2}\), gives
\eqref{eq:H-equation}.  At \((z,t)=(0,0)\), one has \(H=0\) and
\(H_z=1\), so the analytic implicit-function theorem selects a unique small
solution for \(t\) near zero.  This solution corresponds to \(x\to+\infty\)
and hence to the outer branch.
\end{proof}

The exact derivative needed for Newton iteration and error control is
\begin{equation}
 H_z(z;t,\theta)=
 \e^z-t^2\e^{-z}+2\rho t\sin(\theta+\rho z).
 \label{eq:Hz}
\end{equation}
The key point is that \(H_z\to1\) uniformly in the real phase \(\theta\) as
\(t\to0\).  Thus the implicit solution exists not merely for each fixed phase,
but uniformly over the phase circle when \(t\) is sufficiently small.

\subsection{Factorization into two conjugate small monomials}

The scalar equation is best solved after exposing its two complex-conjugate
factors.  From \eqref{eq:F-square-exp},
\begin{equation}
 Y^2=\e^{2ax}
 \bigl(1-\e^{(-2a+\ii\pi)x}\bigr)
 \bigl(1-\e^{(-2a-\ii\pi)x}\bigr).
 \label{eq:factorized-inverse-equation}
\end{equation}
Define
\begin{equation}
 q=\e^{-2ax}=t\e^{-z},
 \qquad
 \lambda=1-\ii\rho,
 \qquad
 \overline\lambda=1+\ii\rho.
 \label{eq:q-lambda}
\end{equation}
Because \(2a\rho=\pi\), the oscillatory factors are
\begin{equation}
 u=q^\lambda=\e^{(-2a+\ii\pi)x},
 \qquad
 v=q^{\overline\lambda}=\e^{(-2a-\ii\pi)x}.
 \label{eq:u-v-physical}
\end{equation}
On the positive real branch, \(v=\overline u\).  Temporarily treating them as
independent analytic variables is the decisive simplification.

Equation \eqref{eq:factorized-inverse-equation} is equivalent to
\begin{equation}
 q=t(1-u)(1-v).
 \label{eq:q-product}
\end{equation}
Let
\begin{equation}
 X=t^\lambda=t\e^{\ii\theta},
 \qquad
 V=t^{\overline\lambda}=t\e^{-\ii\theta}.
 \label{eq:X-V}
\end{equation}
Raising \eqref{eq:q-product} to the powers \(\lambda\) and
\(\overline\lambda\) gives the coupled system
\begin{equation}
 \boxed{
 \begin{aligned}
 u&=X(1-u)^\lambda(1-v)^\lambda,\\
 v&=V(1-u)^{\overline\lambda}(1-v)^{\overline\lambda}.
 \end{aligned}}
 \label{eq:Good-system}
\end{equation}
Moreover, since \(q/t=\e^{-z}\),
\begin{equation}
 \boxed{
 z=-\log\bigl((1-u)(1-v)\bigr).}
 \label{eq:z-log-product}
\end{equation}
Equations \eqref{eq:Good-system}--\eqref{eq:z-log-product} are the exact
Lagrange--Good form of the inverse problem.

\begin{methodbox}[title={The normalization pattern}]
The example illustrates a reusable sequence of steps.
\begin{enumerate}
\item Separate the dominant exponential envelope and invert it exactly.
\item Measure the inverse displacement by a dimensionless variable \(z\).
\item Factor every trigonometric perturbation into conjugate exponentials.
\item Promote the conjugate small monomials to independent variables.
\item Apply multivariate Lagrange inversion; impose conjugacy only at the end.
\end{enumerate}
This converts a nonlinear phase shift such as
\(\cos(\theta+\rho z)\) into an algebraic multivariate fixed-point system.
\end{methodbox}

\subsection{An exact scalar fixed-point identity}

Expanding the logarithms in \eqref{eq:z-log-product}, while retaining
\(q=t\e^{-z}\), gives another useful exact representation:
\begin{align}
 z
 &=\sum_{m\ge1}\frac{u^m+v^m}{m} \notag\\
 &=2\sum_{m\ge1}\frac{t^m}{m}\e^{-mz}
 \cos\bigl(m\theta+m\rho z\bigr).
 \label{eq:scalar-fixed-point}
\end{align}
For the physical positive branch, \(0<q<1\), so the logarithmic series is
absolutely convergent.  Equation \eqref{eq:scalar-fixed-point} is convenient
for numerical fixed-point iteration and for intuitive power counting.  The
multivariate form \eqref{eq:Good-system}, however, is what produces the closed
coefficient formula.

\section{The log-periodic transseries scale}
\label{sec:scale}

\subsection{Complex powers and periodic coefficient blocks}

For nonnegative integers \(j,k\),
\begin{equation}
 X^jV^k
 =t^{j\lambda+k\overline\lambda}
 =t^{j+k}\e^{\ii(j-k)\theta}.
 \label{eq:monomial-block}
\end{equation}
Thus total degree \(n=j+k\) fixes the algebraic size \(t^n=Y^{-2n}\), while
\(j-k\) determines the Fourier harmonic.  This leads to the following
terminology.

\begin{definition}[Log-periodic power transseries]
\label{def:log-periodic}
A log-periodic power transseries in \(t\to0^+\) is an expression
\begin{equation}
 S(t)=\sum_{n\ge n_0}t^nP_n(\log t),
 \label{eq:logperiodic-definition}
\end{equation}
where each \(P_n\) is a finite trigonometric polynomial in \(\log t\).  It may
equivalently be written as a locally finite series of complex powers
\(\sum c_\alpha t^\alpha\), with exponents occurring in conjugate pairs.
\end{definition}

In the present normalization,
\(\theta=-(\pi/(2a))\log t=-\rho\log t\).  Hence a Fourier factor
\(\e^{\ii m\theta}\) is the complex power \(t^{-\ii m\rho}\).  The natural
monomial semigroup is
\begin{equation}
 \Mcal=\left\{t^{j\lambda+k\overline\lambda}:j,k\in\N\right\}.
 \label{eq:monomial-semigroup}
\end{equation}
Every nonconstant monomial has positive real exponent \(j+k\), so only
finitely many monomials lie above any prescribed asymptotic size.  This local
finiteness is what makes multiplication, composition with analytic germs, and
coefficientwise reversion well defined.

\begin{proposition}[Block structure]
\label{prop:block-structure}
Suppose
\[
 z(X,V)=\sum_{j+k\ge1}d_{jk}X^jV^k.
\]
On the physical locus \(V=\overline X\), it has the block decomposition
\begin{equation}
 z(t,\theta)=\sum_{n\ge1}t^nZ_n(\theta),
 \qquad
 Z_n(\theta)=\sum_{j=0}^nd_{j,n-j}\e^{\ii(2j-n)\theta}.
 \label{eq:block-decomposition}
\end{equation}
If \(d_{k j}=\overline{d_{jk}}\), then every \(Z_n\) is real-valued.  Its
Fourier support is
\begin{equation}
 \{-n,-n+2,\ldots,n-2,n\}.
 \label{eq:Fourier-support}
\end{equation}
Consequently
\begin{equation}
 Z_n(\theta+\pi)=(-1)^nZ_n(\theta).
 \label{eq:block-parity}
\end{equation}
\end{proposition}

\begin{proof}
Substitute \eqref{eq:monomial-block} and group terms with \(j+k=n\).  Complex
conjugation interchanges \(j\) and \(k\), proving reality under the stated
symmetry.  Since \(2j-n\) runs from \(-n\) to \(n\) in steps of two, the
support and parity follow.
\end{proof}

\subsection{Asymptotic meaning of truncation}

Because \(\abs{X}=\abs{V}=t\), truncating by total degree is the same as
truncating by asymptotic order.  If the bivariate series is analytic in a
polydisc \(\abs{X},\abs{V}<R\), then for every fixed \(r<R\), Cauchy estimates
give a constant \(C_r\) such that
\begin{equation}
 \left|z(t,\theta)-\sum_{n=1}^{N}t^nZ_n(\theta)\right|
 \le C_r t^{N+1}
 \label{eq:uniform-block-remainder-preview}
\end{equation}
for \(0\le t\le r' < r\), uniformly in \(\theta\).  After dividing by \(2a\),
the inverse remainder is \(O(Y^{-2N-2})=O(y^{-2N-2})\).

This is stronger than a purely formal asymptotic statement: the transseries
will be shown to converge for sufficiently large \(y\).  Nevertheless the
formal viewpoint remains useful because the same coefficient formulas survive
in wider settings where convergence may be unavailable or where additional
asymptotic scales are present.

\section{Closed coefficients from Lagrange--Good inversion}
\label{sec:Good}

\subsection{The coefficient theorem}

For \(w\in\C\) and \(m\in\N\), use the generalized binomial coefficient
\begin{equation}
 \binom{w}{m}=\frac{w(w-1)\cdots(w-m+1)}{m!},
 \qquad \binom{w}{0}=1.
 \label{eq:generalized-binomial}
\end{equation}
The multivariate inversion formula used below goes back to Good
\cite{Good1960}; see also \cite{Gessel1987,Gessel2016,Sokal2009} for modern
proofs and variants.

\begin{theorem}[Closed coefficient formula]
\label{thm:closed-coefficients}
Let \((u,v)\) be the analytic solution near \((X,V)=(0,0)\) of
\eqref{eq:Good-system}, and let \(z\) be given by
\eqref{eq:z-log-product}.  Then
\begin{equation}
 z(X,V)=\sum_{\substack{j,k\ge0\\j+k\ge1}}d_{jk}X^jV^k,
 \label{eq:z-double-series}
\end{equation}
where
\begin{equation}
 d_{jk}=\frac{(-1)^{j+k+1}}{A_{jk}}
 \binom{A_{jk}}{j}\binom{A_{jk}}{k},
 \qquad
 A_{jk}=j\lambda+k\overline\lambda.
 \label{eq:djk-theorem}
\end{equation}
The series is convergent in a neighborhood of \((X,V)=(0,0)\).
\end{theorem}

\begin{proof}
The right-hand sides of \eqref{eq:Good-system} are analytic near the origin
when the principal germs of the powers are chosen.  The Jacobian of the system
with respect to \((u,v)\) is the identity at the origin, so the analytic
implicit-function theorem produces a unique analytic solution with
\(u(0,0)=v(0,0)=0\).

Write
\begin{equation}
 \phi_1(u,v)=(1-u)^\lambda(1-v)^\lambda,
 \qquad
 \phi_2(u,v)=(1-u)^{\overline\lambda}
                         (1-v)^{\overline\lambda}.
 \label{eq:phi12}
\end{equation}
The two-variable Lagrange--Good formula says that, for any analytic germ
\(K(u,v)\),
\begin{equation}
 [X^jV^k]K(u(X,V),v(X,V))
 =[u^jv^k]K(u,v)\phi_1(u,v)^j\phi_2(u,v)^k\Delta(u,v),
 \label{eq:Good-formula}
\end{equation}
where
\begin{equation}
 \Delta=
 \det\left(\delta_{rs}-w_s\frac{\partial}{\partial w_s}
                  \log\phi_r\right)_{r,s=1}^2,
 \qquad (w_1,w_2)=(u,v).
 \label{eq:Good-determinant}
\end{equation}
For \(K=z=-\log((1-u)(1-v))\), one has
\begin{equation}
 z_u=\frac1{1-u},\qquad z_v=\frac1{1-v}.
 \label{eq:z-derivatives}
\end{equation}
Since \(\log\phi_1=-\lambda z\) and
\(\log\phi_2=-\overline\lambda z\), the determinant is rank one and reduces
to
\begin{equation}
 \Delta=1+\lambda u z_u+\overline\lambda v z_v.
 \label{eq:Delta-rank-one}
\end{equation}
For fixed \((j,k)\ne(0,0)\), put
\begin{equation}
 A=A_{jk}=j\lambda+k\overline\lambda,
 \qquad
 P=\phi_1^j\phi_2^k
 =(1-u)^A(1-v)^A=\e^{-Az}.
 \label{eq:A-P}
\end{equation}
The following elementary derivative identity is the cancellation at the heart
of the computation:
\begin{equation}
 u\frac{\partial}{\partial u}
 \left[\left(z+\frac1A\right)P\right]
 =-Azu z_uP,
 \label{eq:integration-by-parts-u}
\end{equation}
and the analogous identity holds with \(u\) replaced by \(v\).  Extracting
\([u^jv^k]\) and using
\([u^jv^k]u\partial_uQ=j[u^jv^k]Q\), we obtain
\begin{align}
 [u^jv^k]zP\Delta
 &=[u^jv^k]zP
   -\frac{\lambda j+\overline\lambda k}{A}
       [u^jv^k]\left(z+\frac1A\right)P \notag\\
 &=-\frac1A[u^jv^k]P.
 \label{eq:key-coefficient-cancellation}
\end{align}
Finally,
\begin{equation}
 [u^jv^k]P
 =(-1)^{j+k}\binom{A}{j}\binom{A}{k}.
 \label{eq:P-coefficient}
\end{equation}
Substitution into \eqref{eq:Good-formula} gives
\eqref{eq:djk-theorem}.
\end{proof}

\begin{remark}[Why the Fibonacci denominator never vanishes]
For \((j,k)\ne(0,0)\),
\[
 \Real A_{jk}=j+k>0.
\]
Hence \(A_{jk}\ne0\).  The universal theorem in
\Cref{sec:general-theorem} also covers zero combined exponents by a removable
continuation, but the Fibonacci specialization never needs that continuation.
\end{remark}

\begin{corollary}[The outer inverse transseries]
\label{cor:inverse-transseries}
On the physical locus \(V=\overline X\), the coefficients satisfy
\begin{equation}
 d_{kj}=\overline{d_{jk}},
 \label{eq:d-conjugacy}
\end{equation}
and \eqref{eq:z-double-series} is real.  Substitution into
\eqref{eq:x-z-substitution} yields \eqref{eq:main-double}, while grouping by
\(n=j+k\) yields \eqref{eq:main-block}--\eqref{eq:Zn-closed}.
\end{corollary}

\begin{proof}
Conjugation interchanges \(\lambda\) and \(\overline\lambda\), hence
interchanges \(A_{jk}\) and \(A_{kj}\); generalized binomial coefficients
respect conjugation.  This proves \eqref{eq:d-conjugacy}.  The remaining
statements follow from \eqref{eq:monomial-block}.
\end{proof}

\subsection{A compact direct derivation of the first terms}

Although the closed formula is the cleanest all-orders result, it is useful to
see the first correction arise directly from \eqref{eq:H-equation}.  Suppose
\[
 z=tZ_1(\theta)+t^2Z_2(\theta)+O(t^3).
\]
At order \(t\),
\[
 Z_1-2\cos\theta=0,
\]
so \(Z_1=2\cos\theta\).  At order \(t^2\), expansion of the exponential and
of the shifted cosine gives
\[
 Z_2+\frac12Z_1^2+1+2\rho Z_1\sin\theta=0.
\]
Using \(Z_1=2\cos\theta\) and elementary double-angle identities yields
\[
 Z_2=-2-\cos(2\theta)-2\rho\sin(2\theta).
\]
This computation already shows how nonlinear phase modulation creates sine
terms even though the original defining equation contains only a cosine.

\section{Explicit coefficient blocks and their structure}
\label{sec:blocks}

\subsection{The first six blocks}

The all-orders formula \eqref{eq:Zn-closed} can be expanded into real
trigonometric polynomials.  The first six are as follows:
\begin{align}
 Z_1(\theta)={}&2\cos\theta,
 \label{eq:Z1}\\[0.25em]
 Z_2(\theta)={}&-2-\cos(2\theta)-2\rho\sin(2\theta),
 \label{eq:Z2}\\[0.25em]
 Z_3(\theta)={}&(6-\rho^2)\cos\theta+5\rho\sin\theta\notag\\
 &+\left(\frac23-3\rho^2\right)\cos(3\theta)
 +3\rho\sin(3\theta),
 \label{eq:Z3}\\[0.25em]
 Z_4(\theta)={}&-9+4(3\rho^2-2)\cos(2\theta)
 +\frac{4\rho(2\rho^2-13)}3\sin(2\theta)\notag\\
 &+\frac{16\rho^2-1}{2}\cos(4\theta)
 +\frac{\rho(16\rho^2-11)}3\sin(4\theta),
 \label{eq:Z4}\\[0.25em]
 Z_5(\theta)={}&
 \frac{\rho^4-95\rho^2+240}{6}\cos\theta
 -\frac{4\rho(2\rho^2-31)}3\sin\theta\notag\\
 &+\frac{27\rho^4-213\rho^2+40}{4}\cos(3\theta)
 -\frac{7\rho(9\rho^2-11)}2\sin(3\theta)\notag\\
 &+\frac{625\rho^4-875\rho^2+24}{60}\cos(5\theta)
 -\frac{25\rho(5\rho^2-1)}6\sin(5\theta),
 \label{eq:Z5}\\[0.25em]
 Z_6(\theta)={}&-\frac{200}{3}
 -\frac{92\rho^4-937\rho^2+450}{6}\cos(2\theta)\notag\\
 &-\frac{\rho(8\rho^4-418\rho^2+1035)}{6}\sin(2\theta)
 -\frac{4(64\rho^4-116\rho^2+9)}3\cos(4\theta)\notag\\
 &-\frac{4\rho(64\rho^4-620\rho^2+261)}{15}\sin(4\theta)
 -\frac{324\rho^4-135\rho^2+2}{6}\cos(6\theta)\notag\\
 &-\frac{\rho(648\rho^4-1530\rho^2+137)}{30}\sin(6\theta).
 \label{eq:Z6}
\end{align}
Numerically,
\begin{equation}
 a=0.481211825059603\ldots,
 \qquad
 \rho=3.26425130263650\ldots.
 \label{eq:numerical-a-rho}
\end{equation}
The coefficients can therefore be sizable even though the accompanying factor
\(t^n=(5y^2)^{-n}\) decays rapidly.

\begin{summarybox}[title={Six-block approximation}]
Define
\[
 \Psi^{[N]}(y)=\frac{\log(\sqrt5y)}a
 +\frac1{2a}\sum_{n=1}^{N}\frac{Z_n(\theta)}{(5y^2)^n},
 \qquad
 \theta=\frac{\pi}{a}\log(\sqrt5y).
\]
Then \(\Psi(y)-\Psi^{[N]}(y)=O(y^{-2N-2})\) as \(y\to\infty\), uniformly
with respect to the phase generated by \(y\).  Equations
\eqref{eq:Z1}--\eqref{eq:Z6} provide a direct approximation through order
\(y^{-12}\), with remainder \(O(y^{-14})\).
\end{summarybox}

\subsection{Fourier support, parity, and polynomial dependence}

The closed coefficient formula makes several structural laws immediate.

\begin{proposition}[Structural laws for \(Z_n\)]
\label{prop:Zn-laws}
For every \(n\ge1\):
\begin{enumerate}
\item \(Z_n\) is a real trigonometric polynomial with the Fourier support
\eqref{eq:Fourier-support}.
\item \(Z_n(\theta+\pi)=(-1)^nZ_n(\theta)\).
\item Written in the real basis \(1,\cos(m\theta),\sin(m\theta)\), every
coefficient is a polynomial in \(\rho\) of degree at most \(n-1\), with
rational coefficients.
\item A constant Fourier term occurs only when \(n=2m\) is even, and then
\begin{equation}
 \avg{Z_{2m}}
 :=\frac1{2\pi}\int_0^{2\pi}Z_{2m}(\theta)\,\dd\theta
 =-\frac1{2m}\binom{2m}{m}^{\!2}.
 \label{eq:mean-Z}
\end{equation}
For odd \(n\), \(\avg{Z_n}=0\).
\end{enumerate}
\end{proposition}

\begin{proof}
The first two statements were proved abstractly in
\Cref{prop:block-structure}.  Since \(A_{jk}=n-\ii\rho(2j-n)\), each
\(d_{jk}\) is a polynomial in \(A_{jk}\) of degree at most \(j+k-1=n-1\):
the apparent denominator is canceled by the factor \(A_{jk}\) contained in
at least one of the two binomial coefficients.  Pairing conjugate terms
produces rational polynomials in \(\rho\).

The zero Fourier mode requires \(2j-n=0\), so \(n=2m\) and \(j=k=m\).  In
that case \(A_{mm}=2m\), and \eqref{eq:djk-theorem} gives
\[
 d_{mm}=-\frac1{2m}\binom{2m}{m}^2.
\]
No other term contributes to the phase average.
\end{proof}

The mean correction has an interesting combinatorial size.  The central
binomial estimate
\(\binom{2m}{m}\sim4^m/\sqrt{\pi m}\) gives
\begin{equation}
 \avg{Z_{2m}}
 \sim-\frac{16^m}{2\pi m^2}.
 \label{eq:mean-growth}
\end{equation}
Thus the even, phase-averaged subseries in \(t\) has radius at most \(1/4\).
This does not determine the exact convergence boundary of the full physical
series, but it gives a sharp structural warning: the convergent expansion is a
large-argument expansion, not a globally convergent representation down to
\(y=1\).

\subsection{The leading oscillation and multiplicative self-similarity}

The first correction is
\begin{equation}
 \Delta_1(y)=\frac{1}{a}t\cos\theta
 =\frac{1}{5ay^2}
 \cos\left(\frac{\pi}{a}\log(\sqrt5y)\right).
 \label{eq:leading-oscillation}
\end{equation}
Multiplication of \(Y\) by \(\phig^2=\e^{2a}\) advances \(\theta\) by
\(2\pi\), while it multiplies \(t\) by \(\phig^{-4}\).  Hence the shape of
the leading oscillation repeats on a geometric progression, with a fixed
amplitude scaling.  More generally, the block \(t^nZ_n(\theta)\) transforms by
\begin{equation}
 t^nZ_n(\theta)\longmapsto
 \phig^{-4n}t^nZ_n(\theta)
 \quad\text{under}\quad Y\longmapsto\phig^2Y.
 \label{eq:discrete-scale-law}
\end{equation}
This is a discrete scale-covariance law, a common signature of complex
asymptotic exponents.

\section{Four complementary algorithms}
\label{sec:algorithms}

The closed formula is ideal for theory, but different tasks favor different
computational forms.  This section records four mutually checking algorithms.

\subsection{Direct coefficient extraction}

\begin{algorithm}[Closed-coefficient construction]
\label{alg:closed}
To construct the expansion through total degree \(N\):
\begin{enumerate}
\item Set \(\lambda=1-\ii\rho\).
\item For every \(n=1,\ldots,N\) and \(j=0,\ldots,n\), put
\[
 k=n-j,
 \qquad A_{jk}=j\lambda+k\overline\lambda,
 \qquad
 d_{jk}=\frac{(-1)^{n+1}}{A_{jk}}
 \binom{A_{jk}}j\binom{A_{jk}}k.
\]
\item Form
\[
 Z_n(\theta)=\sum_{j=0}^nd_{j,n-j}\e^{\ii(2j-n)\theta}.
\]
\item Pair the \(j\) and \(n-j\) terms to convert the answer to a manifestly
real sine--cosine form.
\end{enumerate}
\end{algorithm}

Only \((N+1)(N+2)/2-1\) coefficient pairs are needed.  For symbolic work, the
integer lower arguments of the binomial coefficients make each coefficient a
finite polynomial computation; no infinite expansion is hidden in
\eqref{eq:djk-theorem}.

\subsection{Triangular block recurrence}

The normalized scalar equation produces a particularly simple recurrence.
Suppose \(Z_1,\ldots,Z_{n-1}\) are known and define
\begin{equation}
 z_{<n}(t,\theta)=\sum_{r=1}^{n-1}t^rZ_r(\theta).
 \label{eq:z-less-n}
\end{equation}

\begin{proposition}[Triangular recurrence]
\label{prop:triangular-recurrence}
For \(n\ge1\),
\begin{equation}
 \boxed{
 Z_n(\theta)=-[t^n]\left
 \{\e^{z_{<n}}+t^2\e^{-z_{<n}}
 -2t\cos(\theta+\rho z_{<n})-1\right\}.}
 \label{eq:Zn-recurrence}
\end{equation}
\end{proposition}

\begin{proof}
Insert
\(z=z_{<n}+t^nZ_n+O(t^{n+1})\) into \eqref{eq:H-equation}.  At order \(t^n\),
the new term contributes \(Z_n\) through \(\e^z\).  Its contribution through
\(t^2\e^{-z}\) starts at order \(t^{n+2}\), and through the shifted cosine at
order \(t^{n+1}\).  Therefore the coefficient of \(t^n\) is \(Z_n\) plus the
coefficient computed from \(z_{<n}\), proving the formula.
\end{proof}

The recurrence is well suited to a computer algebra system because it never
requires solving a nonlinear equation at a given order.  It also proves by
induction that each block is a finite trigonometric polynomial.

\subsection{Fixed-point iteration}

Equation \eqref{eq:scalar-fixed-point} defines
\begin{equation}
 \Phi_{t,\theta}(z)=
 2\sum_{m\ge1}\frac{t^m}{m}\e^{-mz}
 \cos(m\theta+m\rho z).
 \label{eq:Phi-map}
\end{equation}
Starting from \(z^{(0)}=0\), one may iterate
\begin{equation}
 z^{(r+1)}=\Phi_{t,\theta}(z^{(r)}).
 \label{eq:fixed-point-iteration}
\end{equation}
For sufficiently small \(t\), this is a contraction on a phase-independent
interval.  For example, if \(\abs z\le R\) and \(te^R<1\), then
\begin{align}
 \abs{\Phi_{t,\theta}(z)}
 &\le -2\log(1-te^R),
 \label{eq:Phi-bound}\\
 \abs{\Phi_{t,\theta}'(z)}
 &\le 2\sqrt{1+\rho^2}\,
       \frac{te^R}{1-te^R}.
 \label{eq:Phi-derivative-bound}
\end{align}
Thus any pair \((t,R)\) for which the right side of
\eqref{eq:Phi-bound} is at most \(R\) and the right side of
\eqref{eq:Phi-derivative-bound} is less than one gives a rigorous contraction
domain.  This method evaluates the exact correction but does not by itself
expose the closed coefficient structure.

\subsection{Newton correction and residual certification}

Given any approximation \(z_0\), one Newton step for the exact scalar equation
is
\begin{equation}
 z_1=z_0-\frac{H(z_0;t,\theta)}{H_z(z_0;t,\theta)},
 \label{eq:Newton-step}
\end{equation}
with \(H_z\) from \eqref{eq:Hz}.  A truncated transseries makes an excellent
initial value: if \(z_0=\sum_{n=1}^Nt^nZ_n\), then
\(H(z_0;t,\theta)=O(t^{N+1})\), and one Newton step normally improves the
formal order to approximately \(O(t^{2N+2})\).

\begin{theorem}[Elementary residual certificate]
\label{thm:residual-certificate}
Let \(z_0\in\R\), \(r>0\), and suppose
\begin{equation}
 H_z(z;t,\theta)\ge m>0
 \quad\text{for every }z\in[z_0-r,z_0+r].
 \label{eq:certificate-derivative}
\end{equation}
If
\begin{equation}
 \abs{H(z_0;t,\theta)}\le mr,
 \label{eq:certificate-residual}
\end{equation}
then \eqref{eq:H-equation} has a unique root \(z_*\) in that interval and
\begin{equation}
 \abs{z_*-z_0}\le
 \frac{\abs{H(z_0;t,\theta)}}{m}.
 \label{eq:certificate-error}
\end{equation}
Consequently the inverse error is at most
\(\abs{H(z_0;t,\theta)}/(2am)\).
\end{theorem}

\begin{proof}
By the mean-value theorem,
\[
 H(z_0+r)\ge H(z_0)+mr\ge0,
 \qquad
 H(z_0-r)\le H(z_0)-mr\le0.
\]
Continuity gives a root, and strict positivity of \(H_z\) gives uniqueness.
Applying the mean-value theorem between \(z_0\) and the root yields
\eqref{eq:certificate-error}.
\end{proof}

This certificate is especially useful in rigorous computation: interval
arithmetic can bound \(H_z\) on a short interval around the truncation, while
the residual is evaluated directly from the original equation.  No estimate
of all omitted coefficients is required.

\section{Analytic convergence and uniform remainders}
\label{sec:convergence}

The word ``transseries'' often suggests a formal or divergent expansion.  In
the present problem the situation is better: the relevant transseries is the
restriction of a convergent bivariate Taylor series.

\begin{theorem}[Convergence for large arguments]
\label{thm:convergence}
There is an \(R>0\) such that the solution \(z(X,V)\) of
\eqref{eq:Good-system}--\eqref{eq:z-log-product} is analytic for
\(\abs X<R\), \(\abs V<R\).  Consequently, for every sufficiently large
positive \(y\), the series \eqref{eq:main-double} and
\eqref{eq:main-block} converge absolutely to the outer inverse \(\Psi(y)\).
The convergence is uniform in the phase \(\theta\) whenever
\(0\le t\le r<R\).
\end{theorem}

\begin{proof}
Analyticity near the origin was established at the start of the proof of
\Cref{thm:closed-coefficients}.  Hence there is a polydisc of analyticity
\(\abs X<R\), \(\abs V<R\), possibly after decreasing \(R\).  On the physical
locus, \(\abs X=\abs V=t\), independently of \(\theta\).  Therefore every
\(t<R\) lies in the polydisc for all real phases, and absolute, locally uniform
convergence of the Taylor series follows.  The analytic solution tends to
zero as \(t\to0\), so by the branch selection in
\Cref{prop:normalized-equation} it is the correction belonging to the outer
inverse.
\end{proof}

\begin{proposition}[A Cauchy-type tail bound]
\label{prop:Cauchy-tail}
Choose \(R_0>0\) such that \(z\) is analytic on a neighborhood of the closed
polydisc \(\abs X,\abs V\le R_0\), and put
\[
 M_{R_0}=\max_{\abs X,\abs V\le R_0}\abs{z(X,V)}.
\]
For \(0<t<R_0\), let \(q=t/R_0\).  Then
\begin{align}
 \left|z(t,\theta)-\sum_{n=1}^Nt^nZ_n(\theta)\right|
 &\le M_{R_0}\sum_{n=N+1}^{\infty}(n+1)q^n \notag\\
 &=M_{R_0}\,
 \frac{q^{N+1}\bigl((N+2)-(N+1)q\bigr)}{(1-q)^2}.
 \label{eq:Cauchy-tail}
\end{align}
The bound is uniform in \(\theta\).
\end{proposition}

\begin{proof}
Cauchy's estimate in two variables gives
\(\abs{d_{jk}}\le M_{R_0}R_0^{-j-k}\).  There are \(n+1\) pairs
\((j,k)\) with \(j+k=n\).  On the physical locus each corresponding monomial
has modulus \(t^n\).  Summing the geometric-arithmetic tail proves the result.
\end{proof}

\begin{corollary}[Poincar\'e remainder]
\label{cor:Poincare-remainder}
For every fixed \(N\ge0\),
\begin{equation}
 \Psi(y)=\frac{\log(\sqrt5y)}a
 +\frac1{2a}\sum_{n=1}^{N}\frac{Z_n(\theta)}{(5y^2)^n}
 +O(y^{-2N-2}),
 \qquad y\to+\infty.
 \label{eq:Poincare-remainder}
\end{equation}
The implied constant can be chosen uniformly over the phase circle.
\end{corollary}

\begin{remark}[The maximal convergence domain]
The analytic implicit-function theorem proves a nonzero convergence domain but
does not locate its nearest singularity.  The maximal bivariate domain is
bounded by the singular variety on which either the fixed-point solution
encounters a branch singularity \(u=1\) or \(v=1\), or the Lagrange--Good
Jacobian determinant vanishes.  Determining which component first meets the
physical curve \((X,V)=(t\e^{\ii\theta},t\e^{-\ii\theta})\) is an interesting
separate problem; see \Cref{sec:future}.
\end{remark}

\section{Numerical validation}
\label{sec:numerics}

The formulas can be checked without relying on the expansion itself: solve
\(\Fcal(x)=y\) numerically on \([2,\infty)\), evaluate the truncated blocks,
and compare.  Let
\begin{equation}
 \Psi^{[N]}(y)=\frac{\log(\sqrt5y)}a
 +\frac1{2a}\sum_{n=1}^{N}\frac{Z_n(\theta)}{(5y^2)^n},
 \label{eq:PsiN}
\end{equation}
with \(\Psi^{[0]}=a^{-1}\log(\sqrt5y)\).  The following table reports absolute
errors.  The displayed values were computed with high-precision arithmetic
from the original defining equation.

\begin{table}[htbp]
\centering
\small
\caption{Absolute errors \(\abs{\Psi^{[N]}(y)-\Psi(y)}\).}
\label{tab:numerical-errors}
\begin{tabular}{@{}rlllll@{}}
\toprule
\(y\) & \(N=0\) & \(N=1\) & \(N=2\) & \(N=4\) & \(N=6\) \\
\midrule
2   & \(1.12696\times10^{-1}\) & \(1.52360\times10^{-2}\) & \(2.96187\times10^{-3}\) & \(3.76916\times10^{-4}\) & \(1.28682\times10^{-4}\) \\
5   & \(1.68278\times10^{-2}\) & \(2.26345\times10^{-4}\) & \(1.85965\times10^{-5}\) & \(2.98868\times10^{-8}\) & \(1.82331\times10^{-11}\) \\
10  & \(5.45185\times10^{-4}\) & \(1.13522\times10^{-5}\) & \(1.54252\times10^{-7}\) & \(3.01709\times10^{-11}\) & \(6.00018\times10^{-15}\) \\
50  & \(1.35073\times10^{-4}\) & \(2.57013\times10^{-8}\) & \(7.18834\times10^{-12}\) & \(8.65363\times10^{-19}\) & \(9.96184\times10^{-25}\) \\
100 & \(3.01016\times10^{-5}\) & \(3.56196\times10^{-9}\) & \(2.97159\times10^{-13}\) & \(4.30138\times10^{-21}\) & \(1.48686\times10^{-28}\) \\
\bottomrule
\end{tabular}
\end{table}

For reference, the corresponding outer inverse values are
\begin{align*}
 \Psi(2)&=3, & \Psi(5)&=5,\\
 \Psi(10)&=6.45779309028821164945\ldots, &
 \Psi(50)&=9.80193485400697171349\ldots,\\
 \Psi(100)&=11.24218977013280389416\ldots.
\end{align*}
The exact equalities \(\Psi(2)=3\) and \(\Psi(5)=5\) are consequences of the
integer interpolation.  More generally,
\begin{equation}
 \boxed{\Psi(F_n)=n\qquad(n\ge2).}
 \label{eq:inverse-Fibonacci-lattice}
\end{equation}
At \(y=2\) the small parameter is only \(t=1/20\), so convergence is slower
and individual errors need not decrease by a fixed ratio.  Already by
\(y=5\), six blocks deliver about eleven correct decimal digits; at larger
arguments the algebraic hierarchy becomes extremely effective.

\begin{remark}[A simple consistency check]
The truncation residual
\[
 R_N(y)=H\!\left(\sum_{n=1}^Nt^nZ_n(\theta);t,\theta\right)
\]
must be \(O(t^{N+1})\).  Expanding symbolically verifies that the coefficients
of \(t,t^2,\ldots,t^N\) vanish identically.  Evaluating \(R_N\) numerically and
applying \Cref{thm:residual-certificate} gives an independent a posteriori
check of every row in \Cref{tab:numerical-errors}.
\end{remark}

\section{A universal inversion theorem for exponential analytic perturbations}
\label{sec:general-theorem}

The preceding derivation is not tied to Fibonacci numbers.  Its essential
features are a dominant exponential envelope and analytic dependence on one or
more exponentially small monomials.  We now formulate and prove the general
result.

\subsection{Set-up and exact reduction}

Let \(r\ge1\), let \(H(\bm u)\) be analytic near
\(\bm u=\bm0\in\C^r\), and assume
\begin{equation}
 H(\bm0)=1.
 \label{eq:H0-general}
\end{equation}
Choose the analytic germ of \(\log H\) that vanishes at the origin.  Consider
\begin{equation}
 Y=C\e^{\sigma x}
 H\!\left(c_1\e^{-\alpha_1x},\ldots,
          c_r\e^{-\alpha_rx}\right),
 \label{eq:general-map}
\end{equation}
where \(C\ne0\), \(\sigma\ne0\), and branches of all logarithms and powers
are fixed.  For a real outer inverse one normally assumes
\begin{equation}
 C>0,
 \qquad \sigma>0,
 \qquad \Real\alpha_j>0.
 \label{eq:real-general-assumptions}
\end{equation}
Define
\begin{equation}
 L=\Log(Y/C),
 \qquad
 \lambda_j=\frac{\alpha_j}{\sigma},
 \qquad
 X_j=c_j\e^{-\lambda_jL}
     =c_j(Y/C)^{-\lambda_j}.
 \label{eq:general-scaled-vars}
\end{equation}
Write
\begin{equation}
 x=\frac{L+z}{\sigma}.
 \label{eq:general-x-z}
\end{equation}
Then the small internal variables are
\begin{equation}
 u_j=c_j\e^{-\alpha_jx}
 =X_j\e^{-\lambda_jz}.
 \label{eq:general-u}
\end{equation}
Substitution into \eqref{eq:general-map} gives the exact relations
\begin{equation}
 z=-\log H(\bm u),
 \qquad
 u_j=X_jH(\bm u)^{\lambda_j}
 \quad(1\le j\le r).
 \label{eq:general-Good-system}
\end{equation}
This is an \(r\)-variable Lagrange--Good system with a common analytic factor.
The common-factor structure is what causes the determinant and the final
coefficient formula to collapse.

For a multi-index
\(\bm n=(n_1,\ldots,n_r)\in\N^r\), write
\begin{equation}
 \abs{\bm n}=n_1+\cdots+n_r,
 \qquad
 \bm X^{\bm n}=\prod_{j=1}^rX_j^{n_j},
 \qquad
 A_{\bm n}=\bm n\mathbin{\bm\cdot}\bm\lambda
 =\sum_{j=1}^rn_j\lambda_j.
 \label{eq:multiindex-A}
\end{equation}
For any scalar \(A\), define the \emph{regularized power quotient}
\begin{equation}
 \mathscr Q_A[H](\bm u)
 :=\sum_{m\ge1}\frac{A^{m-1}}{m!}\bigl(\log H(\bm u)\bigr)^m
 =\begin{cases}
 \dfrac{H(\bm u)^A-1}{A},&A\ne0,\\[0.7em]
 \log H(\bm u),&A=0.
 \end{cases}
 \label{eq:regularized-power-quotient}
\end{equation}
Because \(\log H\) has zero constant term, only finitely many summands in
\eqref{eq:regularized-power-quotient} contribute to any fixed coefficient.
In particular, coefficient extraction from \(\mathscr Q_A[H]\) is an entire,
in fact polynomial, function of \(A\).

\begin{theorem}[Universal coefficient formula]
\label{thm:universal-inversion}
Assume \(H\) is analytic near the origin and \(H(\bm0)=1\).  Then the system
\eqref{eq:general-Good-system} has a unique analytic solution near
\(\bm X=\bm0\), and
\begin{equation}
 z(\bm X)=\sum_{\bm n\ne\bm0}b_{\bm n}\bm X^{\bm n},
 \label{eq:z-general-series}
\end{equation}
with the explicit, everywhere regular coefficients
\begin{equation}
 \boxed{
 b_{\bm n}=-[\bm u^{\bm n}]\mathscr Q_{A_{\bm n}}[H](\bm u).}
 \label{eq:universal-bn}
\end{equation}
Equivalently,
\begin{equation}
 b_{\bm n}=
 \begin{cases}
 -\dfrac1{A_{\bm n}}[\bm u^{\bm n}]H(\bm u)^{A_{\bm n}},
       &A_{\bm n}\ne0,\\[0.9em]
 -[\bm u^{\bm n}]\log H(\bm u),&A_{\bm n}=0.
 \end{cases}
 \label{eq:universal-bn-cases}
\end{equation}
Consequently the inverse of \eqref{eq:general-map} is
\begin{equation}
 \boxed{
 x(Y)=\frac1\sigma\Log\frac YC
 -\frac1\sigma\sum_{\bm n\ne\bm0}
 [\bm u^{\bm n}]\mathscr Q_{A_{\bm n}}[H](\bm u)
 \prod_{j=1}^r
 \left(c_j(Y/C)^{-\alpha_j/\sigma}\right)^{n_j}.}
 \label{eq:universal-inverse}
\end{equation}
The series converges whenever the vector \(\bm X(Y)\) lies in a sufficiently
small polydisc.  Under \eqref{eq:real-general-assumptions}, this holds for all
sufficiently large positive \(Y\); moreover \(\Real A_{\bm n}>0\) for every
nonzero \(\bm n\), so the first line of \eqref{eq:universal-bn-cases} always
applies.
\end{theorem}

\begin{proof}
The analytic implicit-function theorem applies to
\(u_j=X_jH(\bm u)^{\lambda_j}\), because the Jacobian with respect to
\(\bm u\) is the identity at \((\bm u,\bm X)=(\bm0,\bm0)\).  It remains to
compute the Taylor coefficients of \(z=-\log H\).

Set
\begin{equation}
 \phi_i(\bm u)=H(\bm u)^{\lambda_i}.
 \label{eq:general-phi}
\end{equation}
The \(r\)-variable Lagrange--Good formula gives
\begin{equation}
 [\bm X^{\bm n}]z(\bm u(\bm X))
 =[\bm u^{\bm n}]z(\bm u)
 \prod_{i=1}^r\phi_i(\bm u)^{n_i}\Delta(\bm u),
 \label{eq:general-Good-coeff}
\end{equation}
where
\begin{equation}
 \Delta=\det\left(
 \delta_{ij}-u_j\frac{\partial}{\partial u_j}\log\phi_i
 \right)_{i,j=1}^r.
 \label{eq:general-Delta}
\end{equation}
Since \(\log\phi_i=\lambda_i\log H=-\lambda_i z\), the matrix in
\eqref{eq:general-Delta} is
\begin{equation}
 I+\bm\lambda\,\bigl(u_1z_{u_1},\ldots,u_rz_{u_r}\bigr).
 \label{eq:rank-one-matrix}
\end{equation}
The matrix determinant lemma therefore yields
\begin{equation}
 \Delta=1+Ez,
 \qquad
 E:=\sum_{j=1}^r\lambda_ju_j\partial_{u_j}.
 \label{eq:general-Delta-rank-one}
\end{equation}
For fixed \(\bm n\ne\bm0\), abbreviate
\(A=A_{\bm n}\) and
\begin{equation}
 P=\prod_{i=1}^r\phi_i^{n_i}=H^A=\e^{-Az}.
 \label{eq:general-P}
\end{equation}

First suppose \(A\ne0\).  For every \(j\),
\begin{equation}
 u_j\partial_{u_j}\left[\left(z+\frac1A\right)P\right]
 =-Azu_jz_{u_j}P.
 \label{eq:general-derivative-cancellation}
\end{equation}
Using
\([\bm u^{\bm n}]u_j\partial_{u_j}Q=n_j[\bm u^{\bm n}]Q\)
and summing with weights \(\lambda_j\), we find
\begin{align}
 [\bm u^{\bm n}]zP\Delta
 &=[\bm u^{\bm n}]zP
 -\frac{\sum_j\lambda_jn_j}{A}
 [\bm u^{\bm n}]\left(z+\frac1A\right)P \notag\\
 &=-\frac1A[\bm u^{\bm n}]P
 =-[\bm u^{\bm n}]\mathscr Q_A[H].
 \label{eq:general-cancellation-result}
\end{align}
The last equality uses \(\bm n\ne\bm0\), so the constant term subtracted in
\((H^A-1)/A\) does not affect the coefficient.

If \(A=0\), then \(P=1\), and the apparent pole disappears directly:
\begin{align}
 [\bm u^{\bm n}]z\Delta
 &=[\bm u^{\bm n}]\left(z+zEz\right)\notag\\
 &=[\bm u^{\bm n}]z
 +\frac12[\bm u^{\bm n}]E(z^2)\notag\\
 &=[\bm u^{\bm n}]z
 +\frac{A}{2}[\bm u^{\bm n}]z^2
 =-[\bm u^{\bm n}]\log H.
 \label{eq:zero-A-cancellation}
\end{align}
This is exactly the \(A=0\) value in
\eqref{eq:regularized-power-quotient}.  Thus
\eqref{eq:universal-bn} holds in every case.  Substitution in
\eqref{eq:general-x-z} proves \eqref{eq:universal-inverse}.  Finally, under
\eqref{eq:real-general-assumptions}, \(\Real\lambda_j>0\), so
\(X_j(Y)\to0\) and \(\Real A_{\bm n}>0\) for every nonzero multi-index.
\end{proof}

\begin{methodbox}[title={One coefficient extraction replaces reversion}]
Once a map is written in the form \eqref{eq:general-map}, the inverse
coefficient indexed by \(\bm n\) is obtained in three operations:
\begin{enumerate}
\item compute the scalar weight \(A_{\bm n}=\bm n\cdot\bm\lambda\);
\item form the regularized quotient \(\mathscr Q_{A_{\bm n}}[H]\);
\item extract the single coefficient
\(-[\bm u^{\bm n}]\mathscr Q_{A_{\bm n}}[H]\).  When
\(A_{\bm n}\ne0\), this is the familiar operation of extracting the
coefficient of \(H^{A_{\bm n}}\) and dividing by \(-A_{\bm n}\); when
\(A_{\bm n}=0\), it is \(-[\bm u^{\bm n}]\log H\).
\end{enumerate}
No previous inverse coefficients are needed.  The triangular recurrence and
Newton method remain useful computationally, but the reversion problem has
been reduced to coefficient extraction in the \emph{forward} perturbation.
\end{methodbox}

\subsection{Product and radical perturbations}

The most explicit case is a product of binomial factors.

\begin{corollary}[Product formula]
\label{cor:product-formula}
Suppose
\begin{equation}
 H(\bm u)=\prod_{j=1}^r(1-u_j)^{\gamma_j},
 \label{eq:H-product}
\end{equation}
with arbitrary fixed complex exponents \(\gamma_j\).  Then
\begin{equation}
 \boxed{
 b_{\bm n}=
 \frac{(-1)^{\abs{\bm n}+1}}{A_{\bm n}}
 \prod_{j=1}^r
 \binom{\gamma_jA_{\bm n}}{n_j}.}
 \label{eq:product-bn}
\end{equation}
If \(A_{\bm n}=0\), the right-hand side is interpreted by its removable
polynomial limit at zero.
\end{corollary}

\begin{proof}
Because
\[
 H(\bm u)^A
 =\prod_{j=1}^r(1-u_j)^{\gamma_jA},
\]
its \(\bm u^{\bm n}\)-coefficient is
\((-1)^{\abs{\bm n}}\prod_j\binom{\gamma_jA}{n_j}\).  For
\(\bm n\ne\bm0\), this product is divisible by \(A\) as a polynomial in
\(A\), so the quotient has a removable value at zero.  Apply
\eqref{eq:universal-bn} and \eqref{eq:regularized-power-quotient}.
\end{proof}

The formula includes radicals without extra work.  For example, a square root
corresponds to \(\gamma_j=1/2\), and an inverse square root to
\(\gamma_j=-1/2\).  The only branch choice is the analytic germ at
\(\bm u=\bm0\), where every factor equals one.

\subsection{A conjugate-pair theorem for damped oscillations}

A broad real family is
\begin{equation}
 Y=C\e^{\sigma x}
 \left(1-2r\e^{-\mu x}\cos(\omega x+\delta)
       +r^2\e^{-2\mu x}\right)^\gamma,
 \label{eq:conjugate-pair-map}
\end{equation}
where \(C,\sigma,\mu>0\), while \(r,\omega,\delta,\gamma\) are real.  Factor
the oscillatory polynomial as
\begin{align}
 &1-2r\e^{-\mu x}\cos(\omega x+\delta)+r^2\e^{-2\mu x}
 \notag\\
 &\qquad=
 \left(1-r\e^{\ii\delta}\e^{-(\mu-\ii\omega)x}\right)
 \left(1-r\e^{-\ii\delta}\e^{-(\mu+\ii\omega)x}\right).
 \label{eq:conjugate-factorization}
\end{align}
Define
\begin{equation}
 \lambda_+=\frac{\mu-\ii\omega}{\sigma},
 \qquad
 \lambda_-=\frac{\mu+\ii\omega}{\sigma},
 \qquad
 X_+=r\e^{\ii\delta}(Y/C)^{-\lambda_+},
 \qquad
 X_-=\overline{X_+}
 \label{eq:conjugate-parameters}
\end{equation}
for positive real \(Y/C\).

\begin{corollary}[Damped-oscillatory inverse]
\label{cor:damped-oscillatory}
For sufficiently large positive \(Y\), the outer inverse of
\eqref{eq:conjugate-pair-map} is
\begin{equation}
 \boxed{
 x(Y)=\frac1\sigma\log\frac YC
 +\frac1\sigma\sum_{\substack{j,k\ge0\\j+k\ge1}}
 \frac{(-1)^{j+k+1}}{A_{jk}}
 \binom{\gamma A_{jk}}j
 \binom{\gamma A_{jk}}k
 X_+^jX_-^k,}
 \label{eq:damped-oscillatory-inverse}
\end{equation}
where
\begin{equation}
 A_{jk}=j\lambda_++k\lambda_-.
 \label{eq:damped-Ajk}
\end{equation}
The series is real because its terms pair under \(j\leftrightarrow k\).
Grouping by \(j+k\) gives powers of
\((Y/C)^{-\mu/\sigma}\) with periodic coefficients in \(\log(Y/C)\).
\end{corollary}

\subsection{Recovery of the Fibonacci formula}

For \(x>0\), equation \eqref{eq:factorized-inverse-equation} may be square
rooted without ambiguity:
\begin{equation}
 Y=\sqrt5\Fcal(x)
 =\e^{ax}
 \left(1-\e^{(-2a+\ii\pi)x}\right)^{1/2}
 \left(1-\e^{(-2a-\ii\pi)x}\right)^{1/2}.
 \label{eq:F-general-form}
\end{equation}
Thus \eqref{eq:general-map} has
\begin{equation}
 C=1,
 \quad \sigma=a,
 \quad
 \alpha_1=2a-\ii\pi,
 \quad
 \alpha_2=2a+\ii\pi,
 \quad
 \gamma_1=\gamma_2=\frac12,
 \quad c_1=c_2=1.
 \label{eq:F-general-parameters}
\end{equation}
The general exponent ratios are
\begin{equation}
 \frac{\alpha_1}{a}=2-\ii\frac\pi a=2\lambda,
 \qquad
 \frac{\alpha_2}{a}=2+\ii\frac\pi a=2\overline\lambda,
 \label{eq:F-general-lambdas}
\end{equation}
where \(\lambda=1-\ii\rho\) is the half-scaled exponent used earlier.  Hence
\(A^{\mathrm{gen}}_{jk}=2A_{jk}\), and the product formula gives
\begin{align}
 b_{jk}
 &=\frac{(-1)^{j+k+1}}{2A_{jk}}
 \binom{A_{jk}}j\binom{A_{jk}}k
 =\frac12d_{jk}.
 \label{eq:F-b-vs-d}
\end{align}
Multiplication by \(1/\sigma=1/a\) recovers the factor
\(d_{jk}/(2a)\) in \eqref{eq:main-double}.  This verifies that the general
universal theorem and the earlier squared-equation derivation are exactly
consistent.

\subsection{Elementary one-variable checks}

A general formula should reproduce cases whose inverse is known without
transseries machinery.

\begin{example}[The inverse of \(\e^x-1\)]
Take \(C=\sigma=\alpha=c=1\) and \(H(u)=1-u\).  Then
\(Y=\e^x-1\), so
\[
 x=\log(1+Y)=\log Y+\sum_{n\ge1}\frac{(-1)^{n+1}}{nY^n}.
\]
The theorem has \(A_n=n\) and
\[
 b_n=-\frac1n[u^n](1-u)^n
 =\frac{(-1)^{n+1}}n,
\]
which is exactly the logarithmic expansion above.
\end{example}

\begin{example}[A binomially perturbed exponential]
For
\[
 Y=\e^x(1-c\e^{-\alpha x})^\gamma,
\]
let \(\lambda=\alpha\).  The inverse correction is
\begin{equation}
 x-\log Y
 =\sum_{n\ge1}\frac{(-1)^{n+1}}{n\lambda}
 \binom{\gamma n\lambda}{n}
 c^nY^{-n\lambda}.
 \label{eq:one-factor-general}
\end{equation}
When \(\gamma=1\) and \(\lambda=1\), this reduces to the previous example.
For nonintegral \(\lambda\), it is a genuine complex- or fractional-power
inverse series.
\end{example}

\section{The associated transseries algebra}
\label{sec:transseries-algebra}

The analytic theorem has a formal counterpart that clarifies which operations
are legitimate before convergence is considered.

\subsection{The exponent semigroup}

Let
\begin{equation}
 \Gamma=\left\{A_{\bm n}=\sum_{j=1}^rn_j\lambda_j:
 \bm n\in\N^r\right\},
 \qquad \Real\lambda_j>0.
 \label{eq:Gamma-semigroup}
\end{equation}
The real part provides a weight
\begin{equation}
 \operatorname{wt}(\bm X^{\bm n})=\Real A_{\bm n}.
 \label{eq:weight}
\end{equation}
For every \(B>0\), only finitely many multi-indices satisfy
\(\Real A_{\bm n}\le B\), because
\[
 \abs{\bm n}\min_j\Real\lambda_j\le \Real A_{\bm n}.
\]
This local finiteness makes the Cauchy product well defined coefficient by
coefficient.

\begin{definition}[Completed exponential monomial algebra]
\label{def:completed-algebra}
Let \(K\) be a commutative \(\mathbb Q\)-algebra.  The
completed monomial algebra
\begin{equation}
 K[[\bm X]]_{\bm\lambda}
 \label{eq:completed-algebra}
\end{equation}
consists of formal sums
\(S=\sum_{\bm n\in\N^r}s_{\bm n}\bm X^{\bm n}\), ordered by the weight
\eqref{eq:weight}.  The variables \(X_j\) are initially independent
bookkeeping variables.  After specialization to
\(X_j=c_j\e^{-\lambda_jL}\), terms having the same combined exponent
\(A_{\bm n}\) are collected.  Addition is termwise, and multiplication is
the locally finite Cauchy product.
\end{definition}

The maximal ideal \(\mathfrak m\) consists of series with zero constant term.
Every unit \(1+S\), \(S\in\mathfrak m\), admits the usual formal operations
\begin{equation}
 \log(1+S),\qquad (1+S)^\gamma,
 \qquad \exp(S),\qquad \frac1{1+S},
 \label{eq:unit-operations}
\end{equation}
formed from their ordinary Taylor series.  Local finiteness guarantees that
each coefficient receives only finitely many contributions.

\subsection{Differentiation and complex powers}

Let \(L=\Log(Y/C)\) and define the logarithmic derivative
\begin{equation}
 D=\frac{\dd}{\dd L}.
 \label{eq:D-log}
\end{equation}
Since \(X_j=c_j\e^{-\lambda_jL}\),
\begin{equation}
 D\bm X^{\bm n}=-A_{\bm n}\bm X^{\bm n}.
 \label{eq:D-monomial}
\end{equation}
Thus the algebra is stable under differentiation.  Integration is also
coefficientwise provided the constant mode is handled separately and
\(A_{\bm n}\ne0\).  The same scalar \(A_{\bm n}\) that appears in the inverse
coefficient formula is therefore the eigenvalue of \(-D\) on the
corresponding monomial.

When \(\lambda_j=\eta_j-\ii\omega_j\),
\begin{equation}
 X_j=c_j(Y/C)^{-\eta_j}\e^{\ii\omega_j\log(Y/C)}.
 \label{eq:complex-power-phase}
\end{equation}
The real parts \(\eta_j\) determine asymptotic size, while the imaginary parts
\(\omega_j\) generate motion on a phase torus.

\begin{proposition}[Periodic versus quasiperiodic blocks]
\label{prop:periodic-quasiperiodic}
Suppose the exponent set is closed under conjugation and the forward map is
real on the positive real axis.
\begin{enumerate}
\item If all nonzero frequencies \(\omega_j\) are integer multiples of one
frequency \(\omega_0\), the coefficient blocks are periodic in
\(\log(Y/C)\) with period \(2\pi/\omega_0\).
\item If the frequencies have rational rank \(s>1\), the coefficient blocks
are quasiperiodic restrictions of trigonometric polynomials on an
\(s\)-torus.
\item Conjugation symmetry of the data implies conjugation symmetry of the
inverse coefficients, and hence reality of the inverse on the physical locus.
\end{enumerate}
\end{proposition}

\begin{proof}
Every monomial phase is an integer linear combination of the basic
frequencies.  The first two statements are the elementary periodicity theory
of such combinations.  For the third, conjugate the fixed-point system and use
uniqueness of its formal solution.
\end{proof}

\subsection{Formal inversion theorem}

\begin{theorem}[Formal version of \Cref{thm:universal-inversion}]
\label{thm:formal-universal}
Let \(H\in1+(u_1,\ldots,u_r)K[[\bm u]]\) and let
\(\lambda_j\in K\).  Define \(\mathscr Q_A[H]\) by the division-free series
in \eqref{eq:regularized-power-quotient}.  Then the system
\[
 u_j=X_jH(\bm u)^{\lambda_j}
\]
has a unique solution \(\bm u\in\mathfrak m^r\), and
\[
 -\log H(\bm u(\bm X))
 =-\sum_{\bm n\ne0}
 [\bm u^{\bm n}]\mathscr Q_{A_{\bm n}}[H](\bm u)\,
 \bm X^{\bm n}.
\]
No invertibility assumption on \(A_{\bm n}\) is required.
\end{theorem}

\begin{proof}
Successive substitution determines every homogeneous part of \(\bm u\)
triangularly, proving existence and uniqueness.  To verify the coefficient
identity without hidden division, fix \(\bm n\) and replace the \(\lambda_j\)
and the finitely many coefficients of \(H\) needed through degree
\(\abs{\bm n}\) by algebraically independent indeterminates.  Over the
fraction field of the resulting universal polynomial ring, the nonzero linear
form \(A_{\bm n}=\sum_jn_j\Lambda_j\) is invertible, so the algebraic
Lagrange--Good proof of \Cref{thm:universal-inversion} applies.  On the other
hand,
\[
 [\bm u^{\bm n}]\mathscr Q_A[H]
 =\sum_{m=1}^{\abs{\bm n}}
 \frac{A^{m-1}}{m!}
 [\bm u^{\bm n}](\log H)^m
\]
is polynomial in all of those indeterminates.  Hence the identity descends from
the fraction field to the universal polynomial ring and may then be
specialized to arbitrary coefficients of \(H\) and arbitrary
\(\lambda_j\in K\), including specializations for which some
\(A_{\bm n}\) vanishes or is not a unit.
\end{proof}

This result separates two issues that are often conflated.  Formal inversion
requires local finiteness, but after regularization it requires no division by
a combined exponent.  Analytic convergence additionally requires bounds on
the forward germ and a domain in which the implicit solution remains regular.

\subsection{Zero combined exponents and failure of the chosen scale}

The factor \(A_{\bm n}^{-1}\) in the convenient nonzero-weight formula looks
like a resonance denominator, but in this common-factor inversion problem its
pole is always removable.  Indeed,
\begin{equation}
 \frac{H^A-1}{A}
 =\log H+\frac{A}{2!}(\log H)^2
 +\frac{A^2}{3!}(\log H)^3+\cdots,
 \label{eq:removable-A-zero}
\end{equation}
so a zero combined exponent contributes
\(-[\bm u^{\bm n}]\log H\), not a logarithmic factor in \(Y\).

What may fail is the chosen \emph{asymptotic ordering}.
\begin{enumerate}
\item A monomial with \(\Real\lambda_j=0\) does not decay.  It normally
belongs in the leading problem rather than in a perturbative tail.
\item If some \(\Real\lambda_j<0\), the corresponding supposed small variable
grows, so the chart is invalid in that direction.
\item Distinct multi-indices may have the same exponent \(A_{\bm n}\) after
one-variable specialization.  Their coefficients must then be collected; if
the common exponent is zero, the result is a constant-scale contribution that
should usually be absorbed into the leading balance.
\end{enumerate}
Logarithmic monomials can certainly arise in broader transseries problems---for
example when a differential equation requires inversion of \(D\) on a
zero-eigenvalue forcing, or when the dominant envelope already contains
powers of \(x\) and \(\log x\).  They are not forced merely by
\(A_{\bm n}=0\) in the analytic common-factor reversion treated here.

\section{Extensions of the method}
\label{sec:extensions}

\subsection{Several damped harmonics}

A finite exponential--trigonometric perturbation can always be written using
complex exponentials.  For example,
\begin{equation}
 Y=C\e^{\sigma x}
 H\!\left(
 r_1\e^{-(\mu_1-\ii\omega_1)x},
 r_1\e^{-(\mu_1+\ii\omega_1)x},\ldots,
 r_s\e^{-(\mu_s+\ii\omega_s)x}
 \right)
 \label{eq:multi-frequency-map}
\end{equation}
falls directly under \Cref{thm:universal-inversion}.  If the
\(\omega_j/\sigma\) are commensurable, one obtains periodic coefficient
blocks.  Otherwise the inverse contains quasiperiodic coefficient functions.
The universal formula remains a finite coefficient extraction at every
multi-index.

The damping rates need not be equal.  In that case total degree is no longer
the most efficient truncation rule.  One should order monomials by
\begin{equation}
 \operatorname{wt}(\bm n)
 =\sum_jn_j\frac{\mu_j}{\sigma},
 \label{eq:nonuniform-weight}
\end{equation}
and retain all terms below a prescribed weight.  Local finiteness still holds
because every \(\mu_j>0\).

\subsection{A polynomial or logarithmic envelope}

The exact closed formula relies on a pure exponential envelope, for which a
constant displacement \(z/\sigma\) multiplies the envelope by \(\e^z\).  A
more general map may have the form
\begin{equation}
 Y=M(x)H\!\left(c_1\e^{-\alpha_1x},\ldots,c_r\e^{-\alpha_rx}\right),
 \label{eq:general-envelope}
\end{equation}
where, for example,
\begin{equation}
 M(x)=C\e^{\sigma x}x^\beta(\log x)^\gamma.
 \label{eq:polylog-envelope}
\end{equation}
The method then proceeds in two layers.

First invert the dominant envelope:
\begin{equation}
 x_0=M^{-1}(Y).
 \label{eq:base-inverse}
\end{equation}
For \eqref{eq:polylog-envelope}, this base inverse is expressed through
Lambert-type or polynomial--logarithmic transseries.  Next write
\(x=x_0+\delta\).  The exact correction equation is
\begin{equation}
 0=\log\frac{M(x_0+\delta)}{M(x_0)}
 +\log H\!\left(c_j\e^{-\alpha_j(x_0+\delta)}\right).
 \label{eq:anchored-general-envelope}
\end{equation}
Its derivative at \(\delta=0\) is asymptotic to
\((\log M)'(x_0)\), which is nonzero in the outer regime.  The analytic
implicit-function theorem and one-variable Lagrange--B\"urmann inversion then
produce \(\delta\) in a coefficient algebra whose elements are slow
transseries in \(x_0\).  The first correction is universally
\begin{equation}
 \delta
 =-\frac{\log H(c_j\e^{-\alpha_jx_0})}
        {(\log M)'(x_0)}
 +\text{higher perturbative orders}.
 \label{eq:first-general-envelope-correction}
\end{equation}
Thus the pure exponential theorem is the constant-slope member of a broader
anchored reversion calculus.

\subsection{Lagrange--B\"urmann over a coefficient algebra}

Suppose an anchored correction can be written as
\begin{equation}
 \delta=\varepsilon\Phi(\delta),
 \label{eq:LB-coefficient-algebra}
\end{equation}
where \(\varepsilon\) is a small monomial and the Taylor coefficients of
\(\Phi\) belong to a differential algebra of slower asymptotic functions.
Then formally
\begin{equation}
 \delta=\sum_{n\ge1}\frac{\varepsilon^n}{n!}
 \left.\frac{\partial^{n-1}}{\partial s^{n-1}}
 \Phi(s)^n\right|_{s=0}.
 \label{eq:LB-general}
\end{equation}
The formula is ordinary Lagrange--B\"urmann inversion, but interpreted over the
slow coefficient algebra.  This observation connects the present
complex-power theory with polynomial--logarithmic transseries: one scale can
serve as the coefficient field for reversion in a smaller scale.

\subsection{Complex branches}

For complex \(Y\), every choice in
\begin{equation}
 L=\Log(Y/C),
 \qquad
 H(\bm u)^{\lambda_j}=\exp(\lambda_j\Log H(\bm u))
 \label{eq:complex-branch-choices}
\end{equation}
selects an asymptotic inverse chart.  The coefficient formula is unchanged;
only the monomial phases and analytic continuation paths differ.  Singular
curves arise where the Lagrange--Good Jacobian vanishes or where the selected
branch of \(H\) becomes singular.  Analytic continuation around such curves
can permute inverse branches.

For the Fibonacci continuation, squaring removes the outer square root but
also merges sign information.  On the positive real branch this is harmless,
because \(\Fcal(x)>0\) for \(x>0\).  In the complex plane, however, one must
specify whether one is inverting the squared map, the two-factor square-root
map, or a chosen complex Binet continuation.  Their Riemann surfaces are
related but not identical.

\subsection{Other linear recurrences}

Let a second-order recurrence have characteristic roots \(r_1,r_2\) with
\(\abs{r_1}>\abs{r_2}\).  A Binet-type interpolation has the form
\begin{equation}
 U(x)=A r_1^x+B r_2^x.
 \label{eq:general-Binet}
\end{equation}
After factoring the dominant term,
\begin{equation}
 U(x)=A r_1^x
 \left(1+\frac BA\e^{-\alpha x}\right),
 \qquad
 \alpha=\Log(r_1/r_2),
 \label{eq:general-Binet-factor}
\end{equation}
so its inverse is a one-factor instance of
\Cref{thm:universal-inversion}.  If a modulus or quadratic norm is taken, a
conjugate pair appears and the inverse becomes log-periodic whenever
\(\alpha\) has nonzero imaginary part.  Higher-order recurrences lead to
several small root ratios and hence to multivariate, often quasiperiodic,
transseries.

\section{Further consequences for the Fibonacci inverse}
\label{sec:consequences}

\subsection{Sharp first-order oscillation}

The first block gives more than an \(O(y^{-2})\) estimate.

\begin{corollary}[No limiting normalized coefficient]
\label{cor:limsup-liminf}
As \(y\to+\infty\),
\begin{equation}
 5ay^2\left(\Psi(y)-\frac1a\log(\sqrt5y)\right)
 =\cos\left(\frac\pi a\log(\sqrt5y)\right)+O(y^{-2}).
 \label{eq:normalized-leading-correction}
\end{equation}
Consequently
\begin{align}
 \limsup_{y\to\infty}
 5ay^2\left(\Psi(y)-\frac1a\log(\sqrt5y)\right)&=1,
 \label{eq:limsup}\\
 \liminf_{y\to\infty}
 5ay^2\left(\Psi(y)-\frac1a\log(\sqrt5y)\right)&=-1.
 \label{eq:liminf}
\end{align}
In particular, the normalized first correction has no limit.
\end{corollary}

\begin{proof}
Equation \eqref{eq:normalized-leading-correction} is
\eqref{eq:main-block} with the first block retained.  The phase
\((\pi/a)\log(\sqrt5y)\) increases continuously to infinity, so there are
sequences of \(y\)'s on which its cosine is respectively \(1\) and \(-1\).
The error tends to zero along both sequences.
\end{proof}

\subsection{Ordinary series along geometric phase rays}

Fix \(Y_0>0\) and define
\begin{equation}
 Y_m=Y_0\phig^{2m},
 \qquad y_m=Y_m/\sqrt5.
 \label{eq:phase-ray}
\end{equation}
Then
\begin{equation}
 \theta_m=\frac\pi a\log Y_m
 =\frac\pi a\log Y_0+2\pi m,
 \label{eq:fixed-phase-ray}
\end{equation}
so every trigonometric block has a fixed coefficient.  Along this geometric
sequence,
\begin{equation}
 \Psi(y_m)=\frac{\log Y_m}{a}
 +\frac1{2a}\sum_{n\ge1}Y_m^{-2n}Z_n(\theta_0),
 \label{eq:ordinary-series-ray}
\end{equation}
which is an ordinary convergent power series in \(Y_m^{-2}\).  Thus the full
log-periodic transseries may be viewed as a continuous family of ordinary
power series indexed by a phase on the circle.

\subsection{A phase-averaged elliptic generating function}

Let
\begin{equation}
 \overline z(t)=\frac1{2\pi}\int_0^{2\pi}z(t,\theta)\,\dd\theta.
 \label{eq:z-average}
\end{equation}
Using \eqref{eq:mean-Z},
\begin{equation}
 \overline z(t)
 =-\frac12\sum_{m\ge1}\frac1m
 \binom{2m}{m}^{\!2}t^{2m}.
 \label{eq:z-average-series}
\end{equation}
Put \(s=16t^2\).  Since
\begin{equation}
 \binom{2m}{m}=4^m\frac{(1/2)_m}{m!},
 \label{eq:central-binomial-pochhammer}
\end{equation}
one obtains
\begin{equation}
 s\frac{\dd\overline z}{\dd s}
 =-\frac12\left({}_2F_1\!\left(\frac12,\frac12;1;s\right)-1\right).
 \label{eq:z-average-hypergeometric}
\end{equation}
With \(K(s)\) denoting the complete elliptic integral of the first kind in the
parameter convention
\begin{equation}
 K(s)=\int_0^{\pi/2}\frac{\dd\phi}{\sqrt{1-s\sin^2\phi}},
 \label{eq:elliptic-K}
\end{equation}
we have
\({}_2F_1(1/2,1/2;1;s)=2K(s)/\pi\), and therefore
\begin{equation}
 \boxed{
 \overline z(t)=
 \int_0^{16t^2}\left(\frac12-\frac{K(s)}\pi\right)\frac{\dd s}{s}.}
 \label{eq:z-average-elliptic}
\end{equation}
The logarithmic singularity of \(K(s)\) at \(s=1\) is consistent with the
radius \(t=1/4\) of the averaged subseries.  This elliptic connection arises
solely from the diagonal Lagrange--Good coefficients.

\subsection{All logarithmic derivatives}

Let
\begin{equation}
 D_y=y\frac{\dd}{\dd y}=\frac{\dd}{\dd\log Y}.
 \label{eq:Dy}
\end{equation}
Because
\(D_y(X^j\overline X^{\,k})=-2A_{jk}X^j\overline X^{\,k}\), termwise
differentiation of the convergent series gives, for \(m\ge1\),
\begin{equation}
 \boxed{
 D_y^m\Psi(y)=\frac{\delta_{m1}}a
 +\frac1{2a}\sum_{j+k\ge1}
 d_{jk}(-2A_{jk})^mX^j\overline X^{\,k}.}
 \label{eq:all-log-derivatives}
\end{equation}
Equivalently, in block form,
\begin{equation}
 D_y\Psi(y)=\frac1a+\frac1{2a}\sum_{n\ge1}t^n
 \left(-2nZ_n(\theta)+\frac\pi aZ_n'(\theta)\right).
 \label{eq:first-log-derivative-block}
\end{equation}
The leading correction is
\begin{equation}
 y\Psi'(y)=\frac1a
 -\frac{2t}{a}\cos\theta
 -\frac{\pi t}{a^2}\sin\theta+O(t^2).
 \label{eq:Psi-prime-leading}
\end{equation}
Higher ordinary derivatives follow by converting powers of \(D_y\) through
Stirling-number identities.  Thus the same coefficient table controls the
inverse and all of its derivatives.

\subsection{Identities on the Fibonacci lattice}

For every integer \(n\ge2\), \(\Psi(F_n)=n\).  Therefore the convergent
transseries supplies the identity
\begin{equation}
 n=\frac1a\log(\sqrt5F_n)
 +\frac1{2a}\sum_{j+k\ge1}d_{jk}
 X_n^j\overline X_n^{\,k},
 \label{eq:n-Fn-identity}
\end{equation}
where
\begin{equation}
 X_n=(\sqrt5F_n)^{-2+\ii\pi/a}.
 \label{eq:Xn}
\end{equation}
For all sufficiently large \(n\), the series is convergent; its truncations
provide rapidly improving corrections to the elementary approximation
\(n\approx a^{-1}\log(\sqrt5F_n)\).  Binet's exact identity
\begin{equation}
 \sqrt5F_n=\phig^n\left(1-(-1)^n\phig^{-2n}\right)
 \label{eq:Binet-lattice}
\end{equation}
shows that \eqref{eq:n-Fn-identity} is ultimately an intricate re-expansion of
a simple logarithm.  The value of the transseries formula is that it remains
valid for every large real \(y\), not only on the integer lattice.

\section{Computer algebra and numerical implementation}
\label{sec:implementation}

This section gives self-contained code for generating the blocks and checking
them numerically.  The code is included in the article rather than required as
an external build dependency.

\subsection{Wolfram Language: triangular block generator}

The following program implements \eqref{eq:Zn-recurrence}.  It is robust for
symbolic generation because each step expands only to the order currently
needed.

\begin{lstlisting}[style=wolframstyle,caption={Generating the trigonometric blocks \(Z_n\).}]
ClearAll[t, th, rho, inverseBlocks];

inverseBlocks[nmax_Integer?Positive] := Module[
  {zser = 0, blocks = {}, n, residual, zn},
  Do[
    residual = Normal @ Series[
      Exp[zser] + t^2 Exp[-zser]
        - 2 t Cos[th + rho zser] - 1,
      {t, 0, n}
    ];
    zn = FullSimplify[
      TrigReduce[-Coefficient[residual, t, n]],
      Assumptions -> Element[{rho, th}, Reals]
    ];
    AppendTo[blocks, zn];
    zser = zser + zn t^n,
    {n, 1, nmax}
  ];
  blocks
];

blocks = inverseBlocks[6];
TableForm[MapIndexed[{First[#2], #1} &, blocks]]
\end{lstlisting}

A direct implementation of the closed formula is shorter:

\begin{lstlisting}[style=wolframstyle,caption={Closed Lagrange--Good coefficients.}]
ClearAll[aa, d, zBlock, rho, th];

aa[j_, k_] := j (1 - I rho) + k (1 + I rho);

d[j_, k_] := (-1)^(j + k + 1)
  Binomial[aa[j, k], j] Binomial[aa[j, k], k] / aa[j, k];

zBlock[n_Integer?Positive] := FullSimplify[
  ComplexExpand[
    Sum[d[j, n - j] Exp[I (2 j - n) th], {j, 0, n}],
    TargetFunctions -> {Re, Im}
  ],
  Assumptions -> Element[{rho, th}, Reals]
];

Table[zBlock[n], {n, 1, 6}]
\end{lstlisting}

\subsection{Wolfram Language: inverse approximation and residual}

\begin{lstlisting}[style=wolframstyle,caption={Evaluating and certifying a truncation.}]
ClearAll[a, rhoValue, fibCont, psiApprox, residual];

a = ArcCsch[2];
rhoValue = Pi/(2 a);

fibCont[x_?NumericQ] := Sqrt[2/5]
  Sqrt[Cosh[2 a x] - Cos[Pi x]];

psiApprox[y_?NumericQ, nmax_Integer?NonNegative] := Module[
  {Y, tt, phase, corr},
  Y = Sqrt[5] y;
  tt = Y^(-2);
  phase = (Pi/a) Log[Y];
  corr = Sum[
    tt^n (zBlock[n] /. {rho -> rhoValue, th -> phase}),
    {n, 1, nmax}
  ];
  Log[Y]/a + corr/(2 a)
];

residual[y_?NumericQ, nmax_Integer?NonNegative] := Module[
  {Y, tt, phase, z0},
  Y = Sqrt[5] y;
  tt = Y^(-2);
  phase = (Pi/a) Log[Y];
  z0 = 2 a (psiApprox[y, nmax] - Log[Y]/a);
  Exp[z0] + tt^2 Exp[-z0]
    - 2 tt Cos[phase + rhoValue z0] - 1
];

(* Independent direct inversion on the outer branch *)
psiDirect[y_?NumericQ] := x /. FindRoot[
  fibCont[x] == y,
  {x, Log[Sqrt[5] y]/a},
  WorkingPrecision -> 80
];
\end{lstlisting}

\subsection{Python/mpmath numerical checker}

\begin{lstlisting}[style=pythonstyle,caption={High-precision numerical comparison with the defining equation.}]
import mpmath as mp

mp.mp.dps = 80
a = mp.asinh(mp.mpf("0.5"))
rho = mp.pi / (2 * a)


def fib_cont(x: mp.mpf) -> mp.mpf:
    """The real Binet-Fibonacci continuation."""
    return mp.sqrt(mp.mpf(2) / 5) * mp.sqrt(
        mp.cosh(2 * a * x) - mp.cos(mp.pi * x)
    )


def generalized_binomial(z: complex, n: int) -> complex:
    """Binomial(z,n) for a nonnegative integer n."""
    value = mp.mpf(1)
    for m in range(n):
        value *= z - m
    return value / mp.factorial(n)


def z_block(n: int, theta: mp.mpf) -> mp.mpf:
    """Evaluate Z_n(theta) from the closed coefficient formula."""
    lam = 1 - 1j * rho
    total = 0j
    for j in range(n + 1):
        k = n - j
        A = j * lam + k * lam.conjugate()
        d = ((-1) ** (n + 1) / A
             * generalized_binomial(A, j)
             * generalized_binomial(A, k))
        total += d * mp.e ** (1j * (2 * j - n) * theta)
    return mp.re(total)


def psi_approx(y: mp.mpf, order: int) -> mp.mpf:
    """N-block inverse transseries approximation."""
    Y = mp.sqrt(5) * y
    t = Y ** (-2)
    theta = (mp.pi / a) * mp.log(Y)
    z = mp.fsum(t ** n * z_block(n, theta)
                for n in range(1, order + 1))
    return mp.log(Y) / a + z / (2 * a)


def psi_direct(y: mp.mpf) -> mp.mpf:
    """Independent high-precision root on the outer positive branch."""
    x0 = mp.log(mp.sqrt(5) * y) / a
    return mp.findroot(lambda x: fib_cont(x) - y,
                       (x0 - mp.mpf("0.2"), x0 + mp.mpf("0.2")))


for y in map(mp.mpf, [2, 5, 10, 50, 100]):
    exact = psi_direct(y)
    errors = [abs(psi_approx(y, n) - exact)
              for n in [0, 1, 2, 4, 6]]
    print(y, exact, errors)
\end{lstlisting}

\subsection{Implementation cautions}

Three small details prevent common errors.
\begin{enumerate}
\item Use the same branch of \(\Log\) in \(t^\lambda\) and in the phase
\(\theta\).  For positive \(y\), the real logarithm is canonical.
\item Group conjugate coefficient pairs before converting to machine
precision.  This suppresses spurious imaginary roundoff.
\item Near the convergence boundary, evaluate a moderate truncation and then
apply one or two Newton steps to the exact equation, rather than summing many
large, canceling blocks.
\end{enumerate}

\section{Research directions and open problems}
\label{sec:future}

The closed coefficient formula resolves the all-orders large-argument
inversion, but it also exposes several natural frontier questions.

\subsection{The exact convergence boundary}

The analytic solution in the independent variables \((X,V)\) can fail only
when a selected binomial branch becomes singular or when the Lagrange--Good
Jacobian vanishes.  For the Fibonacci system,
\begin{equation}
 \Delta(u,v)=1+\frac{\lambda u}{1-u}
                +\frac{\overline\lambda v}{1-v}.
 \label{eq:F-Delta-explicit}
\end{equation}
A complete convergence analysis therefore begins with the singular system
\begin{equation}
 \begin{cases}
 u=X(1-u)^\lambda(1-v)^\lambda,\\
 v=V(1-u)^{\overline\lambda}(1-v)^{\overline\lambda},\\
 \Delta(u,v)=0,
 \end{cases}
 \label{eq:singular-system}
\end{equation}
together with the branch divisors \(u=1\) and \(v=1\).  The problem is to
locate the singularity of smallest modulus in each physical direction
\((X,V)=(t\e^{\ii\theta},t\e^{-\ii\theta})\), and then minimize over
\(\theta\).  This would replace the qualitative phrase ``sufficiently large''
by a sharp convergence threshold.

\begin{conjecture}[Generic fold singularities on fixed phase rays]
\label{conj:generic-fold}
Fix a real phase \(\theta\) and restrict the bivariate solution to
\((X,V)=(s\e^{\ii\theta},s\e^{-\ii\theta})\).  Except possibly for a discrete
set of phases, the dominant singularities of the resulting one-variable
function \(z_\theta(s)\) are simple folds: they satisfy
\eqref{eq:singular-system}, avoid \(u=1\) and \(v=1\), and have a nonzero
second derivative in the null direction.  Consequently the large-order
coefficients of \(z_\theta\) are finite sums of square-root singularity
contributions of the form
\(C R^{-n}n^{-3/2}\), possibly with oscillatory factors from conjugate
singularities.
\end{conjecture}

This is the generic behavior predicted by analytic implicit-function theory,
but a proof requires global control of the singular variety.  Multivariate
analytic-combinatorics methods may be useful for the homogeneous coefficient
sums; see \cite{FlajoletSedgewick2009} for the univariate background.

\subsection{Large-order asymptotics of the Fourier blocks}

The exact formula \eqref{eq:Zn-closed} reduces the problem to a finite sum of
complex generalized binomial coefficients.  Several regimes are of interest:
\begin{enumerate}
\item fixed Fourier mode \(m=2j-n\) as \(n\to\infty\);
\item bulk scaling \(j/n\to\xi\in(0,1)\);
\item edge scaling \(j\) fixed or \(n-j\) fixed;
\item maximum norms \(\max_\theta\abs{Z_n(\theta)}\).
\end{enumerate}
Stirling analysis of the gamma-function representation
\begin{equation}
 \binom{A}{j}=\frac{\Gamma(A+1)}
 {\Gamma(j+1)\Gamma(A-j+1)}
 \label{eq:binomial-gamma}
\end{equation}
should identify the dominant Fourier modes and connect coefficient growth to
the singular geometry of \eqref{eq:singular-system}.

\subsection{Uniform approximations near the convergence boundary}

Near a fold singularity, summing the Taylor transseries is not the best
representation.  One should introduce a local square-root variable and build
a uniform approximant that matches the large-argument series away from the
fold.  Rational, Pad\'e, Hermite--Pad\'e, or conformally mapped approximants
could then extend accurate evaluation toward \(y=1\).  The exact residual
certificate in \Cref{thm:residual-certificate} provides a convenient way to
validate such accelerations.

\subsection{All real and complex inverse branches}

The outer branches are the only branches that survive for large real values,
but the bounded oscillatory region contains additional local inverses.  At a
nondegenerate critical point, these are Puiseux series in
\((y-y_c)^{1/2}\); \Cref{app:finite-branches} records the leading terms.  A
complete real branch atlas would locate every critical point and critical
value.  The complex problem is richer: branch points of the square root in
\(\Fcal\), critical points of the complexified map, and logarithmic monodromy
of the inverse all interact.

\subsection{Other recurrences and exponential sums}

Theorem \ref{thm:universal-inversion} applies immediately to Binet formulas
for Lucas sequences and, after choosing a dominant root, to higher-order linear
recurrences.  A systematic classification could answer:
\begin{enumerate}
\item which root configurations yield ordinary powers, log-periodic blocks,
or quasiperiodic blocks;
\item how repeated characteristic roots introduce polynomial factors and
therefore nested polynomial--logarithmic scales;
\item how moduli, norms, and real-part projections alter branch structure;
\item whether arithmetic sampling at integer arguments forces additional
identities among the coefficient blocks.
\end{enumerate}

\subsection{Formalization}

A proof-assistant development could be organized in modular layers:
\begin{enumerate}
\item formal multivariate power series and generalized binomial powers;
\item the Lagrange--Good coefficient identity;
\item the rank-one determinant lemma and cancellation
\eqref{eq:general-cancellation-result};
\item analytic implicit-function and convergence transfer;
\item the Fibonacci specialization and explicit first blocks;
\item interval-certified residual bounds for numerical examples.
\end{enumerate}
The universal theorem is especially suitable for formalization because the
hard-looking inversion collapses to a short algebraic identity once the
standard multivariate inversion theorem is available.

\subsection{Interaction with broader transseries scales}

The exponential monomial algebra here can be nested inside slower scales, as
in \Cref{sec:extensions}.  Conversely, the coefficient blocks may themselves
carry \(q\)-series, special-function, or recursively defined asymptotics.
Developing a systematic tower
\begin{equation}
 \text{polynomial--logarithmic scale}
 \supset
 \text{exponential small scale}
 \supset
 \text{log-periodic or quasiperiodic blocks}
 \label{eq:scale-tower}
\end{equation}
would connect the present note with the wider algebraic theory of transseries
\cite{vanderHoeven2006,Costin2009}.  The essential design principle is always
the same: choose an ordered, locally finite monomial family and perform
reversion over the coefficient algebra generated by slower scales.

\section{Conclusion}
\label{sec:conclusion}

The real continuation
\[
 \Fcal(x)=\sqrt{\frac25}
 \sqrt{\cosh(2x\log\phig)-\cos(\pi x)}
\]
looks elementary, but its inverse combines three distinct mechanisms:
exponential growth, exponentially damped oscillation, and nonlinear phase
feedback.  The dominant envelope gives the logarithm
\(a^{-1}\log(\sqrt5y)\).  The damped oscillation survives inversion as a
log-periodic hierarchy in \((5y^2)^{-1}\).  Nonlinear feedback creates higher
harmonics and sine components.

The decisive exact reduction is
\[
 u=X(1-u)^\lambda(1-v)^\lambda,
 \qquad
 v=V(1-u)^{\overline\lambda}(1-v)^{\overline\lambda},
 \qquad
 z=-\log((1-u)(1-v)).
\]
Multivariate Lagrange--Good inversion then yields the explicit coefficient
\[
 d_{jk}=\frac{(-1)^{j+k+1}}{j\lambda+k\overline\lambda}
 \binom{j\lambda+k\overline\lambda}{j}
 \binom{j\lambda+k\overline\lambda}{k}.
\]
This formula proves the all-orders expansion, its reality, Fourier support,
parity, and convergence for large arguments.  It also gives a practical
algorithm and explains the unexpectedly rapid numerical accuracy.

More importantly, the example reveals a universal principle.  For every map
\[
 Y=C\e^{\sigma x}H(c_j\e^{-\alpha_jx}),
\]
the inverse coefficient indexed by \(\bm n\) is simply
\[
 -\frac1{A_{\bm n}}[\bm u^{\bm n}]H(\bm u)^{A_{\bm n}},
 \qquad A_{\bm n}=\sum_jn_j\frac{\alpha_j}{\sigma}.
\]
The entire reversion is encoded in ordinary coefficients of scalar powers of
the forward perturbation.  Complex-conjugate exponent ratios produce
log-periodic blocks; several independent imaginary parts produce
quasiperiodicity; positive real parts guarantee a decaying, locally finite
small scale.  Even outside that regime, the apparent zero-weight poles in the
coefficient formula are removable.

Thus the inverse Fibonacci continuation is not an isolated special-function
calculation.  It is a model example of a broad and explicitly solvable class
of exponential--oscillatory inverse problems.

\appendix

\section{Local inverse branches at finite critical values}
\label{app:finite-branches}

For completeness, we record the generic local behavior of the additional
inverse branches that occur in the bounded oscillatory region.  This is
separate from the large-argument transseries.

\begin{proposition}[Puiseux inversion at a simple critical point]
\label{prop:Puiseux-critical}
Let \(x_c\ne0\) satisfy
\begin{equation}
 \Fcal'(x_c)=0,
 \qquad
 \Fcal''(x_c)\ne0,
 \qquad
 y_c=\Fcal(x_c).
 \label{eq:critical-assumptions}
\end{equation}
Then the two local inverse branches have the Puiseux expansions
\begin{equation}
 x=x_c\pm
 \sqrt{\frac{2(y-y_c)}{\Fcal''(x_c)}}
 -\frac{\Fcal'''(x_c)}{3\Fcal''(x_c)^2}(y-y_c)
 +O\bigl((y-y_c)^{3/2}\bigr),
 \label{eq:critical-Puiseux}
\end{equation}
with the square root interpreted in the appropriate real or complex sector.
\end{proposition}

\begin{proof}
Write \(h=x-x_c\) and \(s=y-y_c\).  Taylor expansion gives
\[
 s=Ah^2+Bh^3+O(h^4),
 \qquad
 A=\frac12\Fcal''(x_c),
 \qquad
 B=\frac16\Fcal'''(x_c).
\]
Seek \(h=c_1s^{1/2}+c_2s+O(s^{3/2})\).  Matching the coefficients of \(s\)
and \(s^{3/2}\) yields
\[
 Ac_1^2=1,
 \qquad
 2Ac_1c_2+Bc_1^3=0.
\]
Thus \(c_1=\pm A^{-1/2}\) and
\(c_2=-B/(2A^2)=-\Fcal'''(x_c)/(3\Fcal''(x_c)^2)\), proving the formula.
\end{proof}

This square-root branching is qualitatively different from the outer inverse
at infinity, whose correction is analytic in the complex small monomials
\(X\) and \(V\).

\section{Coefficient table in complex form}
\label{app:coefficient-table}

For low degrees, it is sometimes more compact to retain the complex
coefficients rather than expand into sines and cosines.  Let
\(A_{jk}=j+k-\ii\rho(j-k)\).  The homogeneous block of degree \(n\) is
\begin{equation}
 Z_n(\theta)=\sum_{j=0}^nd_{j,n-j}\e^{\ii(2j-n)\theta},
 \qquad
 d_{jk}=\frac{(-1)^{j+k+1}}{A_{jk}}
 \binom{A_{jk}}j\binom{A_{jk}}k.
 \label{eq:appendix-d-table-master}
\end{equation}
The boundary coefficients simplify to
\begin{equation}
 d_{n0}=\frac{(-1)^{n+1}}{n\lambda}\binom{n\lambda}{n},
 \qquad
 d_{0n}=\overline{d_{n0}},
 \label{eq:boundary-coefficients}
\end{equation}
and, for even degree \(2m\), the diagonal coefficient is
\begin{equation}
 d_{mm}=-\frac1{2m}\binom{2m}{m}^2.
 \label{eq:diagonal-coefficient}
\end{equation}
These two families control respectively the highest Fourier harmonics and the
phase averages.

\section{Notation summary}
\label{app:notation}

\begin{longtable}{@{}>{\raggedright\arraybackslash}p{0.20\textwidth}
                    >{\raggedright\arraybackslash}p{0.72\textwidth}@{}}
\toprule
Symbol & Meaning \\
\midrule
\endfirsthead
\toprule
Symbol & Meaning \\
\midrule
\endhead
\(\phig\) & Golden ratio \((1+\sqrt5)/2\). \\
\(a\) & \(\log\phig=\arsinh(1/2)=\arccsch 2\). \\
\(\Fcal(x)\) & Real nonnegative Fibonacci continuation in the problem statement. \\
\(\Psi(y)\) & Positive outer inverse, valued in \([2,\infty)\). \\
\(Y\) & Scaled argument \(\sqrt5y\). \\
\(t\) & Small real parameter \(Y^{-2}=(5y^2)^{-1}\). \\
\(\theta\) & Logarithmic phase \((\pi/a)\log Y\). \\
\(\rho\) & Phase-feedback constant \(\pi/(2a)\). \\
\(z\) & Dimensionless inverse displacement defined by
\(\Psi=a^{-1}\log Y+z/(2a)\). \\
\(\lambda\) & Complex exponent \(1-\ii\rho\). \\
\(X,V\) & Independent small monomials \(t^\lambda,t^{\overline\lambda}\); on
the physical locus \(V=\overline X\). \\
\(A_{jk}\) & Combined exponent \(j\lambda+k\overline\lambda\). \\
\(d_{jk}\) & Closed double-series coefficient in \eqref{eq:djk-theorem}. \\
\(Z_n(\theta)\) & Real trigonometric coefficient block of order \(t^n\). \\
\(H(z;t,\theta)\) & Exact scalar residual in \eqref{eq:H-equation}. \\
\(H(\bm u)\) & General analytic perturbation germ in
\Cref{sec:general-theorem}; context distinguishes it from the scalar residual. \\
\(A_{\bm n}\) & General combined exponent \(\sum_jn_j\alpha_j/\sigma\). \\
\(\mathscr Q_A[H]\) & Regularized power quotient
\((H^A-1)/A\), continued at \(A=0\) by \(\log H\). \\
\(b_{\bm n}\) & Universal inverse coefficient
\(-[\bm u^{\bm n}]\mathscr Q_{A_{\bm n}}[H]\); for
\(A_{\bm n}\ne0\), this is
\(-A_{\bm n}^{-1}[\bm u^{\bm n}]H^{A_{\bm n}}\). \\
\bottomrule
\end{longtable}

\section{A compact derivation checklist}
\label{app:checklist}

For a new exponential--oscillatory inverse problem, the following checklist
usually leads directly to the answer.
\begin{enumerate}
\item Factor the forward map as \(C\e^{\sigma x}H(\bm u)\) with
\(u_j=c_j\e^{-\alpha_jx}\) and \(H(0)=1\).
\item Choose \(L=\Log(Y/C)\), \(x=(L+z)/\sigma\), and
\(X_j=c_j\e^{-(\alpha_j/\sigma)L}\).
\item Verify that \(\Real(\alpha_j/\sigma)>0\) in the asymptotic sector.
\item Form \(A_{\bm n}=\sum_jn_j\alpha_j/\sigma\).
\item Extract
\[
 b_{\bm n}=-[\bm u^{\bm n}]\mathscr Q_{A_{\bm n}}[H].
\]
For \(A_{\bm n}\ne0\), this is
\(-A_{\bm n}^{-1}[\bm u^{\bm n}]H^{A_{\bm n}}\); for
\(A_{\bm n}=0\), use \(-[\bm u^{\bm n}]\log H\).
\item Pair conjugate indices and group equal real weights into periodic or
quasiperiodic blocks.
\item Check the first blocks by substitution into the exact normalized
equation.
\item For numerical use, evaluate a moderate truncation, compute its exact
residual, and finish with Newton iteration plus a derivative bound.
\end{enumerate}

\begin{thebibliography}{99}

\bibitem{Costin2009}
O.~Costin,
\emph{Asymptotics and Borel Summability},
Chapman \& Hall/CRC Monographs and Surveys in Pure and Applied Mathematics,
vol.~141, CRC Press, 2009.

\bibitem{FlajoletSedgewick2009}
P.~Flajolet and R.~Sedgewick,
\emph{Analytic Combinatorics},
Cambridge University Press, 2009.

\bibitem{Gessel1987}
I.~M. Gessel,
A combinatorial proof of the multivariable Lagrange inversion formula,
\emph{Journal of Combinatorial Theory, Series A} \textbf{45} (1987),
178--195.

\bibitem{Gessel2016}
I.~M. Gessel,
Lagrange inversion,
\emph{Journal of Combinatorial Theory, Series A} \textbf{144} (2016),
212--249.

\bibitem{Good1960}
I.~J. Good,
Generalizations to several variables of Lagrange's expansion, with applications
to stochastic processes,
\emph{Proceedings of the Cambridge Philosophical Society} \textbf{56} (1960),
367--380.

\bibitem{OzvatanPashaev2017}
M.~\"Ozvatan and O.~K. Pashaev,
Generalized Fibonacci sequences and Binet--Fibonacci curves,
\emph{arXiv:1707.09151} (2017).

\bibitem{ProveItSeriesTransseries}
V.~Reshetnikov,
\emph{ProveIt: Series and Transseries Drafts},
\url{https://github.com/VladimirReshetnikov/ProveIt/tree/main/Analysis/FabiusFunction/docs/semi-formalized-research-frontiers/drafts/series-and-transseries},
accessed 3 September 2026.

\bibitem{Sokal2009}
A.~D. Sokal,
A ridiculously simple and explicit implicit function theorem,
\emph{S\'eminaire Lotharingien de Combinatoire} \textbf{61A} (2009),
Article B61Ad, 21 pp.

\bibitem{vanderHoeven2006}
J.~van der Hoeven,
\emph{Transseries and Real Differential Algebra},
Lecture Notes in Mathematics, vol.~1888, Springer, 2006.

\end{thebibliography}

\end{document}
